% !TeX root = ../../Book2.tex
% 纲要标题：## 第20章*　分次代数、滤过与非交换实例
% 纲要位置：计划纲要.md，第 1696 行。

\OptionalChapter{分次代数、滤过与非交换实例}{b2:ch:20}
\BookChapterNumber{b2:ch:20}

\begin{chapterabstract}
本章建立滤过、伴随分次和 Rees 代数，证明若干有限性提升结果，并计算 Weyl 代数与量子平面的正规单项式基。交换子引出 Lie 代数及其包络代数；Ore 扩张与 Diamond 引理分别处理扭曲乘法和重写。最后的例子表明，代数有有限展示，也未必在指定字序下有有限 Gröbner 基。
\end{chapterabstract}

\begin{learninggoals}
\begin{itemize}
\item 构造伴随分次与 Rees 代数，说明穷尽性和滤过下界如何保证降次论证。
\item 在任意特征下证明 Weyl 代数的正规单项式基，比较特征零的单性与正特征的中心。
\item 使用扭曲乘法计算量子平面，辨认参数的阶如何影响中心。
\item 用交换子和矩阵核验 Lie 公理，构造伴随表示，并把 Lie 表示对应为泛包络代数上的模。
\item 从算子模型构造 Ore 扩张，按左模约定解释微分方程。
\item 检查字序、包含及重叠歧义，证明正规式唯一，并区分约化终止与 Gröbner 基的有限性。
\end{itemize}
\end{learninggoals}

在第一 Weyl 代数里，关系 \(\partial x=x\partial+1\) 允许把 \(\partial x^2\) 改写为 \(x^2\partial+2x\)。按字长取 \(\deg x=\deg\partial=1\)，两边主项的次数都是三；特征不为二时，修正项 \(2x\) 的次数是一。在相应的最高次商片中，低次项成为零，于是原来不交换的两个生成元的主符号交换。这里舍去的是低次乘法信息，原代数的哪些性质仍可从这些商片恢复，需要另作证明。

\ProseParagraphBreak

也可以直接规定所有 \(x\) 排在 \(\partial\) 之前，反复使用关系重排。字长和逆序数的下降能保证过程停止，并使每个元素成为 \(x^i\partial^j\) 的线性组合；这还没有证明表达式唯一。若不同重排顺序强迫出额外关系，预期的单项式仍可能线性相关。后面的坏 PBW 关系正会使一个原本写出的生成元变成零。

\ProseParagraphBreak

证明独立性需要比重排更多的信息。Weyl 代数把 \(\partial\) 实现为 \(\partial_X+Y\)，使 \(x^i\partial^j\) 作用于一得到独立单项式 \(X^iY^j\)；量子平面直接在预定基上构造扭曲乘法，Ore 扩张则在自由左模上构造满足关系的算子。一般泛包络代数由 PBW 定理确定伴随分次为对称代数，再把独立性提升回来。Diamond 引理使用另一种检验：给出下降字序，并证明所有包含和重叠歧义的两支最终相同。它们所提供的，都是不能仅从生成元重排中得到的独立性依据。

\ProseParagraphBreak

滤过和包络代数的构造使用第八章的张量代数、对称代数及按关系取商；有限性提升还用到第十二章的 Noether 模概念。

\ProseParagraphBreak
Weyl 代数及其表示可参阅 Etingof 等人的《Introduction to Representation Theory》\cite[\S 2.7, pp. 17--18]{EtingofEtAl2011RepresentationTheory}；该处也提醒正特征中的自然多项式表示不忠实。本章的量子平面不把生成元取逆，与该书同节的量子 Laurent 型例子范围不同。字序、交换单项式序与含分母的局部化适用的计算对象也不同。

% !TeX root = ../../Book2.tex
% 纲要标题：### 20.1 分次与滤过
% 纲要位置：计划纲要.md，第 1698 行。

\section{分次与滤过}
\label{b2:sec:2001}

给多项式记录全次数，可以在乘法中保留最高次部分。非交换关系若还含有较低次项，直接分次可能失效，但次数不超过给定界的子空间仍然有用。这就是滤过的出发点。以下固定域 \(k\)，代数均结合含幺；本节的滤过均递增、穷尽，并在负次数为零。

% 纲要要点（供目录核对）：
%
% * 分次模和分次代数；区分恰为一次与次数至多一，非齐次交换关系自然给出滤过
% * 伴随分次与 Rees 构造；主符号用 symb，明确乘积降次、内部理想与环境求值核、非零参数求值的次数权重和平坦性范围
% * 滤过上的有限性问题；匹配主符号后逐层降次，子模用诱导滤过；整环命题的逆向反例在正文证明，原题改为粗化滤过的 Rees 展示与纤维
%

\subsection{分次模和分次代数}
\label{b2:subsec:2001:01}

\cref{b2:def:associative-graded-algebra}已经定义非负分次代数 \(A=\bigoplus_{n\ge0}A_n\)，其乘法满足 \(A_iA_j\subseteq A_{i+j}\)。一个左 \(A\) 模的分次是直和分解 \(M=\bigoplus_{n\in\mathbb Z}M_n\)，使 \(A_iM_j\subseteq M_{i+j}\)；带有这种分解的模称为\term[fenci-mo]{分次模}[graded module]，保持次数的同态称为分次同态。每个元素都只有有限多个非零齐次分量。

\begin{definition}\label{b2:def:filtered-algebra-module}
设 \(k\) 为域，\(A\) 为含幺结合 \(k\) 代数。若一列 \(k\) 子空间
\[
0=F_{-1}A\subseteq F_0A\subseteq F_1A\subseteq\cdots,
\quad \bigcup_{n\ge0}F_nA=A,
\]
满足 \(1\in F_0A\) 及 \(F_iA\,F_jA\subseteq F_{i+j}A\)，则称 \((F_nA)_{n\geq0}\) 为 \(A\) 的非负\term[lvguo]{滤过}[filtration]；负指标的滤过项约定为零。若 \(M\) 为左 \(A\) 模，\((F_nM)_{n\in\mathbb Z}\) 为递增的 \(k\) 子空间列，满足 \(\bigcup_nF_nM=M\)、\(F_nM=0\) 对 \(n<0\)，以及 \(F_iA\,F_jM\subseteq F_{i+j}M\)，则称它为与 \(A\) 的滤过相容的模滤过。
\end{definition}

对有统一整数下界的模滤过，平移模的指标便可使用本节的论证。这里的穷尽性保证每个元素最终出现，下界则保证降次归纳能到达零；两项条件各有用途。

\ProseParagraphBreak
分次给出滤过 \(F_nA=\bigoplus_{i\le n}A_i\)。分次中的 \(A_n\) 记录恰为 \(n\) 次的齐次分量，滤过中的 \(F_nA\) 则记录次数不超过 \(n\) 的全部元素，因此滤过项彼此包含，而不是彼此直和。反过来，滤过不一定来自这样一个与乘法相容的直和。例如关系 \(\mathrm{d}x-x\mathrm{d}=1\) 混合二次与零次，但按字长不超过 \(n\) 所生成的子空间仍构成滤过。后面将严格构造这个例子。

\ProseParagraphBreak
若 \(M\) 由 \(m_1,\ldots,m_r\) 生成，并给每个生成元指定非负次数 \(d_i\)，可以取
\[
F_nM=\sum_i(F_{n-d_i}A)m_i.
\]
这确实相容且穷尽。这样的表达把有限生成元及其所需次数一并记录下来，而不仅记录一个抽象生成集合。

\subsection{伴随分次与 Rees 构造}
\label{b2:subsec:2001:02}

\begin{definition}\label{b2:def:associated-graded-rees}
设 \(k\) 为域，\(A\) 为带有非负穷尽滤过 \((F_nA)\) 的含幺结合 \(k\) 代数。定义 \(A\) 的\term[bansui-fenci-daishu]{伴随分次代数}[associated graded algebra]为
\[
\operatorname{gr}_F A=\bigoplus_{n\ge0}F_nA/F_{n-1}A,
\]
乘法由 \((a+F_{i-1}A)(b+F_{j-1}A)=ab+F_{i+j-1}A\) 给出。若左 \(A\) 模 \(M\) 带有相容滤过 \((F_nM)\)，则定义 \(\operatorname{gr}_F M=\bigoplus_{n\geq0}F_nM/F_{n-1}M\)。若非零 \(a\in A\) 首次属于 \(F_nA\)，则称 \(n\) 为其滤过次数，称 \(\operatorname{symb}(a)=a+F_{n-1}A\) 为其\term[zhufuhao]{主符号}[principal symbol]。
\end{definition}

改变代表元所产生的乘积误差落在 \(F_{i+j-1}A\)，故上述乘法良定义；结合律由原代数继承，模作用的验证相同。若非零 \(a,b\) 的滤过次数分别为 \(i,j\)，则
\[
\operatorname{symb}(a)\operatorname{symb}(b)
=ab+F_{i+j-1}A\in F_{i+j}A/F_{i+j-1}A.
\]
当 \(ab\ne0\) 且次数为 \(i+j\) 时，这才等于 \(\operatorname{symb}(ab)\)；若 \(ab=0\)，或其实际次数小于 \(i+j\)，上述乘积为零。因此原代数没有零因子，也不能保证主符号相乘不为零。

\ProseParagraphBreak
令 \(z\) 为与 \(A\) 中全部元素交换的不定元。在中心多项式环 \(A[z]\) 中定义\term[Rees-daishu]{Rees 代数}[Rees algebra]
\[
\mathcal R_F(A)=\bigoplus_{n\ge0}(F_nA)z^n.
\]
滤过的乘法条件保证它为子代数。由于 \(1\in F_0A\subseteq F_1A\)，\(z\) 也属于其中且居中。
\symidx{\(\operatorname{gr}_F A\)、\(\mathcal R_F(A)\)：伴随分次代数与 Rees 代数}

求零纤维时，\((z)\) 指 Rees 代数内部的理想 \(z\mathcal R_F(A)\)。它的第 \(n\) 次分量是 \((F_{n-1}A)z^n\)：要在这个子代数内部提出一个 \(z\)，剩下的系数必须已属于较低的滤过项。环境多项式环中的求值 \(z=0\) 限制到 Rees 代数后，核却是 \(\mathcal R_F(A)\cap zA[z]\)，包含全部正次项。总有
\[
z\mathcal R_F(A)\subseteq\mathcal R_F(A)\cap zA[z],
\]
但这个包含可以严格，因而不能用环境多项式环中的求值代替内部取商。

\ProseParagraphBreak
例如 \(A=k[x]\) 配通常次数滤过，则
\[
\mathcal R_F(A)=k[z,xz]\cong k[z,w],\qquad w\longmapsto xz.
\]
这里 \(z^iw^j\) 的像为 \(x^jz^{i+j}\)，不同的 \((i,j)\) 给出不同单项式，所以所写同态确为同构。元素 \(w=xz\) 不属于 Rees 代数内部的 \((z)\)：若 \(xz=zq\) 且 \(q\in\mathcal R_F(A)\)，在 \(k[x,z]\) 中消去 \(z\) 就有 \(q=x\)，但 \(x\notin\mathcal R_F(A)\)，因为其常数次系数不在 \(F_0A=k\) 中。因此模掉 \(z\) 后，\(w\) 仍保留下来，所得商为 \(k[w]\)，正好记录 \(x\) 的最高次信息。下面分别计算参数为零和非零时的纤维；平坦性另用到第七章的\cref{b2:thm:pid-flat-torsionfree}。

\begin{theorem}\label{b2:thm:rees-fibres}
设 \(k\) 为域，\(A\) 为带有非负穷尽滤过 \((F_nA)\) 的含幺结合 \(k\) 代数，令 \(z\) 为与 \(A\) 交换的不定元，\(\mathcal R_F(A)=\bigoplus_{n\geq0}(F_nA)z^n\subseteq A[z]\) 为 Rees 代数，则有自然 \(k\) 代数同构
\[
\mathcal R_F(A)/(z)\cong\operatorname{gr}_F A,
\qquad\mathcal R_F(A)/(z-\lambda)\cong A\quad(\lambda\in k^\times).
\]
此外，\(\mathcal R_F(A)\) 作为 \(k[z]\) 模平坦。
\end{theorem}
\begin{proof}
将 \(a_nz^n\) 送到 \(a_n+F_{n-1}A\)，得到到伴随分次的满同态。其核的第 \(n\) 次分量为 \((F_{n-1}A)z^n\)，正是 \(z\mathcal R_F(A)\)，所以得到第一个同构。

在 \(z=1\) 求值得到到 \(A\) 的同态；穷尽性保证它满射。若 \(p(z)=\sum_{n=0}^N a_nz^n\) 的像为零，则 \(\sum_n a_n=0\)。令
\[
q(z)=-\sum_{n=0}^{N-1}\left(\sum_{i=0}^n a_i\right)z^n.
\]
第 \(n\) 个系数属于 \(F_nA\)，所以 \(q\in\mathcal R_F(A)\)，而逐项相减给出 \(p=(z-1)q\)。反向显然，故求值同态的核恰为 \((z-1)\)。

对任意 \(\lambda\in k^\times\)，逐次乘以 \(\lambda^n\) 给出 \(\mathcal R_F(A)\) 的自同构
\[
\theta_\lambda\colon\sum_n a_nz^n\longmapsto\sum_n\lambda^na_nz^n.
\]
它与 \(z=1\) 处的求值复合，得到 \(z=\lambda\) 处的求值 \(\sum_n a_nz^n\mapsto\sum_n\lambda^na_n\)。因此该同态满射，且其核为
\[
\theta_\lambda^{-1}\bigl((z-1)\mathcal R_F(A)\bigr)
= (z-\lambda)\mathcal R_F(A),
\]
得到每个非零纤维的同构。

\(A[z]\) 作为 \(k[z]\) 模无挠：一个非零标量多项式与非零向量系数多项式相乘，其最高次系数仍非零。子模 \(\mathcal R_F(A)\) 因而无挠。\(k[z]\) 是主理想整环，由\cref{b2:thm:pid-flat-torsionfree}得到平坦性，无需假设 Rees 代数有限生成。
\end{proof}

平坦性是 \(k[z]\) 模的性质，并不保证各个纤维作为代数同构，也不保证其乘法相同。这里的非零纤维都恢复 \(A\)，零纤维却只保留相应滤过次数上的乘积；\cref{b2:prop:filtered-domain-converse-failure}将给出原代数没有零因子而零纤维含非零幂零元的例子。

\subsection{滤过上的有限性问题}
\label{b2:subsec:2001:03}

\begin{theorem}\label{b2:thm:filtered-generator-lifting}
设 \(k\) 为域，\(A\) 为带有非负穷尽滤过 \((F_nA)\) 的含幺结合 \(k\) 代数，\(M\) 为带有相容、穷尽且负次数为零的滤过 \((F_nM)\) 的左 \(A\) 模。若 \(\operatorname{gr}_F M\) 由有限个次数为 \(d_i\) 的齐次元素 \(\bar m_i\) 作为 \(\operatorname{gr}_F A\) 模生成，则任意代表 \(m_i\in F_{d_i}M\) 生成 \(M\)，并有
\[
F_nM=\sum_i(F_{n-d_i}A)m_i.
\]
\end{theorem}
\begin{proof}
右侧包含于左侧来自滤过相容性。反向要把分次中的表达提升到原模：先匹配最高的一层，再相减，使剩余部分降到较低层。对 \(n\) 归纳，给定 \(m\in F_nM\)，其在 \(F_nM/F_{n-1}M\) 中的类可写为 \(\sum_i\bar a_i\bar m_i\)，其中取齐次分量后可令 \(\deg\bar a_i=n-d_i\)，负次数项为零。选择 \(a_i\in F_{n-d_i}A\) 提升这些系数，便有
\[
m-\sum_i a_im_i\in F_{n-1}M.
\]
对差值用归纳假设；所得较低次系数仍包含在 \(F_{n-d_i}A\) 中。\(n=0\) 时差值在 \(F_{-1}M=0\)，归纳开始。穷尽性再说明这些元素生成全部 \(M\)。
\end{proof}

\begin{theorem}\label{b2:thm:filtered-noetherian-transfer}
设 \(k\) 为域，\(A\) 为带有非负穷尽滤过 \((F_nA)\) 的含幺结合 \(k\) 代数，\(M\) 为带有相容、穷尽且负次数为零的滤过 \((F_nM)\) 的左 \(A\) 模。若 \(\operatorname{gr}_F A\) 是左 Noether 环，则 \(A\) 也是左 Noether 环。若 \(\operatorname{gr}_F M\) 是 Noether 左 \(\operatorname{gr}_F A\) 模，则 \(M\) 是 Noether 左 \(A\) 模。将左模及左 Noether 条件同时换为右侧条件，结论也成立。
\end{theorem}
\begin{proof}
对任意子模 \(N\le M\)，取诱导滤过 \(F_nN=N\cap F_nM\)。这样，\(F_nN\) 映入 \(F_nM/F_{n-1}M\) 的核恰为 \(N\cap F_{n-1}M=F_{n-1}N\)，所以每个商片单射，\(\operatorname{gr}_F N\) 才能成为 \(\operatorname{gr}_F M\) 的齐次子模；任意另选的滤过未必有这一性质。由 Noether 假设它有限生成。可以取齐次生成元：将一个有限生成组各自拆成齐次分量，所得有限族仍在该齐次子模内并生成它。由\cref{b2:thm:filtered-generator-lifting}，\(N\) 有限生成。每个子模都有限生成，故 \(M\) Noether。取 \(M=A\) 为左正则模即得到环的断言。对于右模，所有系数写在生成元右侧，同一个降次归纳与诱导滤过论证成立。
\end{proof}

\begin{proposition}\label{b2:prop:filtered-domain}
设 \(k\) 为域，\(A\) 为带有非负穷尽滤过 \((F_nA)\) 的含幺结合 \(k\) 代数。若伴随分次代数 \(\operatorname{gr}_F A\) 没有零因子，则 \(A\) 没有零因子，且任意非零 \(a,b\in A\) 的乘积滤过次数等于两者次数之和。
\end{proposition}
\begin{proof}
非零 \(a,b\) 都有最低的滤过次数，主符号非零。它们的符号乘积不为零，故 \(ab\) 的相应商类不为零。这同时说明 \(ab\ne0\) 及乘积不降次。
\end{proof}

穷尽性和下界不能略去。若在非零模上取全部 \(F_nM=0\)，则伴随分次为零，却看不到原模；这违反穷尽性。若允许所有整数指标并令 \(F_nM=M\)，则滤过虽穷尽，伴随分次仍为零，因为没有开始降次归纳的下界。本节后面使用的字长滤过都满足已经列出的条件。

\begin{proposition}\label{b2:prop:filtered-domain-converse-failure}
设 \(k\) 为任意域，给含幺 \(k\) 代数 \(A=k[x]\) 赋予滤过
\[
F_nA=\{f\in k[x]:\deg f\le 2n\}\quad(n\ge0),
\qquad F_nA=0\quad(n<0),
\]
其中零多项式属于每个非负滤过项。则 \(A\) 是整环，而
\[
\operatorname{gr}_F A\cong k[u,v]/(u^2),
\qquad \deg u=\deg v=1.
\]
因此\cref{b2:prop:filtered-domain}的逆命题不成立。
\end{proposition}
\begin{proof}
多项式次数的乘法公式保证 \(F_iA\,F_jA\subseteq F_{i+j}A\)，而这个滤过显然穷尽且 \(1\in F_0A\)。令 \(u=\operatorname{symb}(x)\)、\(v=\operatorname{symb}(x^2)\)。两者都在一次分量中；但 \(x^2\in F_1A\)，所以把 \(u\) 与自身相乘，所得在 \(F_2A/F_1A\) 中的类为零。于是有分次代数同态 \(k[u,v]/(u^2)\to\operatorname{gr}_F A\)。

零次分量以一为基。对每个 \(n\ge1\)，\(F_nA/F_{n-1}A\) 的一组基是 \(x^{2n-1}\)、\(x^{2n}\) 的类，它们分别是 \(uv^{n-1}\)、\(v^n\) 的像。另一方面，\(k[u,v]/(u^2)\) 的第 \(n\) 次分量恰以 \(uv^{n-1},v^n\) 为基，所以此同态逐次为同构。特别地，\(u\ne0\) 而 \(u^2=0\)，尽管原代数 \(k[x]\) 没有零因子。
\end{proof}

% !TeX root = ../../Book2.tex
% 纲要标题：### 20.2 Weyl 代数
% 纲要位置：计划纲要.md，第 1704 行。

\section{Weyl 代数}
\label{b2:sec:2002}

在多项式上，先乘不定元再求导，与先求导再乘不定元，恰好相差恒等算子。Weyl 代数只保留这条交换关系。我们先证明抽象代数中的正规单项式确实独立，再判断它对通常多项式空间的作用是否忠实；这两项问题在正特征中有不同答案。

% 纲要要点（供目录核对）：
%
% * 微分与乘法的交换关系；区分任意特征下的抽象 Weyl 代数、通常多项式作用与辅助变量忠实作用
% * 正规单项式及作用；下降重排给生成性，辅助算子在一上的取值记录两个指数而给独立性
% * 特征零与正特征的差别；中心秩 p²、有限维表示迹障碍和中心零纤维的矩阵商在正文完整证明；原题改为商模分解、平移纤维及中心分式除代数
%

\subsection{微分与乘法的交换关系}
\label{b2:subsec:2002:01}

在 \(k[X]\) 上令 \(X\) 表示乘以不定元的算子，令 \(D_0\) 为形式求导。Leibniz 公式\footnote{莱布尼茨（Gottfried Wilhelm Leibniz，1646-1716），德国数学家及哲学家，创立微积分，并研究符号演算。}给出 \(D_0X-XD_0=1\)，这在任意特征下都成立。

\begin{definition}\label{b2:def:first-weyl-algebra}
设 \(k\) 为域。将自由结合代数 \(k\langle x,\mathrm{d}\rangle\) 对关系 \(\mathrm{d}\cdot x-x\cdot\mathrm{d}=1\) 生成的双边理想取商，所得代数称为 \(k\) 上的第一\term[Weyl-daishu]{Weyl 代数}[first Weyl algebra]，记为
\[
A_1(k)=k\langle x,\mathrm{d}\rangle/(\mathrm{d}\cdot x-x\cdot\mathrm{d}-1).
\]
其代数同态由生成元的像决定，并要求这些像满足 \(\mathrm{d}\cdot x-x\cdot\mathrm{d}=1\)。
\end{definition}
\symidx{\(A_1(k)\)：第一 Weyl 代数}

由\cref{b2:prop:algebra-relations-universal}，\(x\mapsto X,\mathrm{d}\mapsto D_0\) 给出 \(A_1(k)\) 在 \(k[X]\) 上的作用。但关系成立只证明作用存在，尚未证明这个作用单射。

\subsection{正规单项式及其作用}
\label{b2:subsec:2002:02}

\begin{theorem}\label{b2:thm:weyl-normal-basis}
设 \(k\) 为域，\(A_1(k)=k\langle x,\mathrm{d}\rangle/(\mathrm{d}\cdot x-x\cdot\mathrm{d}-1)\)，则全部 \(x^i\mathrm{d}^j\)、\(i,j\ge0\)，构成 \(A_1(k)\) 的一组 \(k\) 基。因此每个元素唯一写为有限和 \(\sum_j f_j(x)\mathrm{d}^j\)，其中 \(f_j\in k[x]\)。
\end{theorem}
\begin{proof}
生成性来自 \(\mathrm{d}\cdot x=x\cdot\mathrm{d}+1\)。对一个字，按字长及其中位于某个 \(x\) 左边的 \(\mathrm{d}\) 的配对数作字典式归纳：将相邻 \(\mathrm{d}\cdot x\) 换成 \(x\cdot\mathrm{d}+1\)，第一项的字长不变、错序数减少一，第二项的字长减少二。这两个量都是非负整数，故不能无限下降，最终只剩全部 \(x\) 在前、全部 \(\mathrm{d}\) 在后的正规字。不过归约终止这里只保证正规字生成，尚不能排除它们之间仍有线性关系。

为证独立性，我们要找到一个作用，使 \(x^i\mathrm{d}^j\) 的像保留两个指数的信息。通常的求导算子在常数一上反复作用都为零，不能用一来检测 \(j\)。因此在两个交换变量的多项式空间 \(k[X,Y]\) 上令 \(X\) 仍表示乘法，并令
\[
D=\frac{\partial}{\partial X}+Y,
\]
其中 \(Y\) 也表示乘法算子。增加的乘 \(Y\) 算子与乘 \(X\) 交换，故仍有 \(DX-XD=1\)，给出一个表示。同时 \(\partial(Y^j)/\partial X=0\)，所以 \(D(Y^j)=Y^{j+1}\)，从常数一出发便有 \(D^j(1)=Y^j\)，因而
\[
X^iD^j(1)=X^iY^j.
\]
若抽象代数中有有限关系 \(\sum_{i,j}c_{ij}x^i\mathrm{d}^j=0\)，令该算子作用在一上，得到 \(\sum_{i,j}c_{ij}X^iY^j=0\)。交换多项式的单项式独立，故所有 \(c_{ij}=0\)。整个论证没有使用阶乘可逆性。
\end{proof}

这个双变量表示因而在任意特征下都忠实。在特征 \(p\) 中，虽然通常的第 \(p\) 次求导在 \(k[X]\) 上为零，这里的 \(D^p(1)=Y^p\) 却不为零。它证明的是抽象 Weyl 代数中正规字的独立性，并不能据此断言通常的多项式作用也忠实；后者将在\cref{b2:prop:weyl-polynomial-action}中单独判断。后面的\cref{b2:thm:diamond-lemma}还会把改写的下降条件与歧义检查结合起来，证明正规式唯一。Weyl 代数的正规单项式基还可参阅 Etingof 等人\cite[Proposition 2.7.1(i), pp. 17--18]{EtingofEtAl2011RepresentationTheory}；那里使用另一种带辅助变量的表示。

\ProseParagraphBreak
用字长滤过，\(F_nA_1(k)\) 恰由 \(i+j\le n\) 的正规单项式生成。交换关系中的修正项降低次数二，故最高次部分变得交换。

\begin{proposition}\label{b2:prop:weyl-associated-graded}
设 \(k\) 为域，\(A_1(k)=k\langle x,\mathrm{d}\rangle/(\mathrm{d}\cdot x-x\cdot\mathrm{d}-1)\)，按 \(x,\mathrm{d}\) 的字长赋予滤过 \(F_nA_1(k)\)，则 \(\operatorname{gr}_F A_1(k)\cong k[X,D]\)。因此 \(A_1(k)\) 没有零因子，且左右 Noether。
\end{proposition}
\begin{proof}
将交换变量 \(X,D\) 分别送到一次主符号 \(\operatorname{symb}(x),\operatorname{symb}(\mathrm{d})\)。它们交换且生成伴随分次，故得到满同态。\cref{b2:thm:weyl-normal-basis}说明第 \(n\) 个商片恰以 \(i+j=n\) 的单项式符号为基，因此同态逐次单射。

交换多项式环 \(k[X,D]\) 是整环；由第一卷定理12.4.10（Hilbert 基定理），它还为 Noether 环。应用\cref{b2:prop:filtered-domain}及\cref{b2:thm:filtered-noetherian-transfer}，得到无零因子性和两侧 Noether 性。
\end{proof}

例如 \(\mathrm{d}^2x^3=x^3\mathrm{d}^2+6x^2\mathrm{d}+6x\)，由两次使用 \(\mathrm{d}\cdot x^n=x^n\mathrm{d}+nx^{n-1}\) 得到。等式在正特征中仍成立，只是整数系数按该特征取值。

\subsection{特征零与正特征}
\label{b2:subsec:2002:03}

为避免与群交换子混淆，本节的 \([a,b]=ab-ba\) 表示代数中的加法交换子。直接归纳给出
\begin{equation}\label{b2:eq:weyl-commutators}
[\mathrm{d},x^i\mathrm{d}^j]=i x^{i-1}\mathrm{d}^j,\qquad
[x,x^i\mathrm{d}^j]=-j x^i\mathrm{d}^{j-1}.
\end{equation}
零指数时右侧解释为零。

\begin{theorem}\label{b2:thm:weyl-center-simplicity}
设 \(k\) 为域，\(A_1(k)=k\langle x,\mathrm{d}\rangle/(\mathrm{d}\cdot x-x\cdot\mathrm{d}-1)\)。若 \(\operatorname{char}k=0\)，则 \(Z\bigl(A_1(k)\bigr)=k\)，且 \(A_1(k)\) 是单环。若 \(\operatorname{char}k=p>0\)，则
\[
Z\bigl(A_1(k)\bigr)=k[x^p,\mathrm{d}^p],
\]
这是二元多项式环；\(A_1(k)\) 是其中心上的自由模，一组基为
\[
\{x^i\mathrm{d}^j:0\le i,j<p\},
\]
故在中心上秩为 \(p^2\)。此时 \(A_1(k)\) 不是单环。
\end{theorem}
\begin{proof}
由\eqref{b2:eq:weyl-commutators}及正规单项式基，元素 \(\sum c_{ij}x^i\mathrm{d}^j\) 同时与 \(x,\mathrm{d}\) 交换，当且仅当每个非零系数满足 \(i c_{ij}=j c_{ij}=0\)。在特征零中只剩 \(i=j=0\)；在特征 \(p\) 中要求 \(p\mid i,j\)。这证明两个中心公式，且正规单项式基保证 \(x^p,\mathrm{d}^p\) 之间没有多项式关系。

在特征零中，要证明非零双边理想 \(I\) 等于整个代数，只需在其中得到非零标量。双边性保证与 \(x\) 或 \(\mathrm{d}\) 取交换子后仍留在 \(I\)，而\eqref{b2:eq:weyl-commutators}表明这分别降低 \(\mathrm{d}\) 或 \(x\) 的指数。选其中非零元素 \(P=\sum_{j=0}^m f_j(x)\mathrm{d}^j\)，使 \(\mathrm{d}\) 次数 \(m\) 最小。若 \(m>0\)，则 \([P,x]\in I\) 的 \(\mathrm{d}\) 次数为 \(m-1\)，最高系数为 \(m\cdot f_m(x)\ne0\)，矛盾。因此 \(I\) 含非零 \(f(x)=a_rx^r+\cdots+a_0\)，其中 \(a_r\ne0\)。记 \(\operatorname{ad}_{\mathrm{d}}(u)=[\mathrm{d},u]\)，则
\[
\operatorname{ad}_{\mathrm{d}}^{\,r}(f)=r!a_r\in k^\times.
\]
这个交换子仍属于 \(I\)，故 \(1\in I\)，\(I=A_1(k)\)。这里 \(r!\ne0\) 正是特征零假设的作用。

在特征 \(p\) 中，将每个指数唯一写为 \(i=pa+r\)、\(j=pb+s\)，其中 \(0\le r,s<p\)。利用 \(x^p,\mathrm{d}^p\) 居中，便有
\[
x^i\mathrm{d}^j=(x^p)^a(\mathrm{d}^p)^b x^r\mathrm{d}^s,
\]
所以这些余数指数的单项式在中心上生成 \(A_1(k)\)。若它们之间有中心系数的线性关系，将各系数展开为 \(x^p,\mathrm{d}^p\) 的多项式，不同商和余数给出不同的正规单项式。\cref{b2:thm:weyl-normal-basis}迫使所有系数为零，证明中心上的自由性及秩公式。中心元 \(x^p\) 生成的双边理想就是 \(x^pA_1(k)\)，按正规单项式基它由 \(x\) 指数至少为 \(p\) 的单项式生成。因此它非零且不含一，是一个真双边理想。
\end{proof}

\begin{proposition}\label{b2:prop:weyl-polynomial-action}
设 \(k\) 为域，\(A_1(k)=k\langle x,\mathrm{d}\rangle/(\mathrm{d}\cdot x-x\cdot\mathrm{d}-1)\)。令 \(x\) 在 \(k[X]\) 上作用为乘 \(X\)，\(\mathrm{d}\) 作用为形式求导，则这个 \(A_1(k)\) 模作用在 \(\operatorname{char}k=0\) 时忠实，在 \(\operatorname{char}k>0\) 时不忠实。
\end{proposition}
\begin{proof}
特征零时，若非零 \(P=\sum_{j=0}^m f_j(x)\mathrm{d}^j\) 作用为零，记 \(M_X\) 为乘 \(X\) 的算子。由\eqref{b2:eq:weyl-commutators}，\([f_j(X)\mathrm{d}^j,M_X]=j f_j(X)\mathrm{d}^{j-1}\)。因此从 \(P\) 开始反复按 \([P,M_X]\) 的顺序取 \(m\) 次交换子，仍作用为零；低于 \(m\) 次的项全部消失，所得算子是乘 \(m!f_m(X)\)，令其作用在一上便得到矛盾。\(m=0\) 时直接如此。

特征 \(p\) 时，第 \(p\) 次形式求导在每个 \(X^n\) 上都为零，因为连续 \(p\) 个整数的乘积含因子 \(p\)。所以非零的抽象元素 \(\mathrm{d}^p\) 作用为零，非零性则由\cref{b2:thm:weyl-normal-basis}保证。
\end{proof}

\begin{proposition}\label{b2:prop:weyl-finite-dimensional-representations}
设 \(k\) 为任意域，\(A=A_1(k)=k\langle x,\mathrm{d}\rangle/(\mathrm{d}\cdot x-x\cdot\mathrm{d}-1)\)。若 \(\operatorname{char}k=0\)，则每个有限维幺性左 \(A\) 模都是零模。若 \(\operatorname{char}k=p>0\)，则每个有限维幺性左 \(A\) 模的 \(k\) 维数都被 \(p\) 整除；此时
\[
V_p=k[X]/(X^p)
\]
在 \(x\) 作用为乘 \(X\)、\(\mathrm{d}\) 作用为形式求导时，是一个 \(p\) 维简单左 \(A\) 模，该作用诱导 \(k\) 代数同构
\[
A/(x^p,\mathrm{d}^p)\cong\operatorname{End}_k(V_p)\cong M_p(k),
\]
其中 \((x^p,\mathrm{d}^p)\) 是这两个中心元生成的双边理想。
\end{proposition}
\begin{proof}
在一个有限维幺性左 \(A\) 模 \(V\) 上，记 \(x,\mathrm{d}\) 的作用为 \(X,D\)，则 \(DX-XD=1_V\)。矩阵迹满足 \(\operatorname{Tr}(DX)=\operatorname{Tr}(XD)\)，故
\[
0=\operatorname{Tr}(DX-XD)=(\dim_kV)\,1_k.
\]
在特征零中这迫使 \(\dim_kV=0\)；在特征 \(p\) 中则迫使 \(p\mid\dim_kV\)。这一迹障碍也见 Etingof 等人\cite[Problem 2.7.4, p. 18]{EtingofEtAl2011RepresentationTheory}。

以下设 \(\operatorname{char}k=p>0\)。形式求导保持理想 \((X^p)\)，因为 \((X^pf)'=X^pf'\)，所以乘法与求导都下降到 \(V_p\)，仍满足 Weyl 关系。若一个非零子模含有非零 \(f\)，取次数小于 \(p\) 的代表，设其次数为 \(r\)、最高系数为 \(a_r\ne0\)。反复求导 \(r\) 次得到非零常数 \(r!a_r\)；这里 \(r<p\)，所以阶乘不为零。因此子模含一，再乘 \(X\) 就得到 \(1,X,\ldots,X^{p-1}\)，故 \(V_p\) 简单。

在 \(V_p\) 上，\(X^p\) 与 \(D^p\) 都为零，故作用经过所写商代数。要证明其像是全部线性算子，可以在像中逐一构造矩阵单位。令 \(H=XD\)，则 \(H(X^r)=rX^r\) 对 \(0\le r<p\)。因而
\[
E_0=\prod_{r=1}^{p-1}\frac{H-r\,1_{V_p}}{-r}
\]
恰为到常数项的投影：它在一上为恒等，在其余基向量上为零。所有分母都是 \(k\) 中的非零元素。对 \(0\le i,j<p\)，令
\[
E_{ij}=\frac{1}{j!}X^iE_0D^j.
\]
当 \(r<j\) 时，\(D^j(X^r)=0\)；当 \(r=j\) 时，它等于 \(j!\)；当 \(r>j\) 时，它是正次数单项式的标量倍，随后被 \(E_0\) 消去。因此 \(E_{ij}(X^r)=\delta_{jr}X^i\)，正是该基下的矩阵单位。它们全在商代数的像中，故诱导映射满射。另一方面，商代数由 \(x^i\mathrm{d}^j\)、\(0\le i,j<p\) 张成，维数至多 \(p^2\)；目标维数恰为 \(p^2\)，所以这个满射也是单射。
\end{proof}

% !TeX root = ../../Book2.tex
% 纲要标题：### 20.3 量子平面与扭曲代数
% 纲要位置：计划纲要.md，第 1710 行。

\section{量子平面与扭曲代数}
\label{b2:sec:2003}

把交换关系中的修正换成一个非零参数，可以得到另一类可计算的代数。下面始终固定 \(q\in k^\times\)。参数为一时恢复交换多项式环；其他参数是否为单位根，将直接影响中心的大小。

% 纲要要点（供目录核对）：
%
% * 量子平面的定义；从移动 jr 对字母产生扭曲因子，用预定基上的结合乘法证明独立性
% * 扭曲多项式环的实例；量子关系齐次而 Weyl 修正低次，取伴随分次时保留的信息不同
% * 中心与参数的关系；有限中心秩要求参数的精确阶，区别于 Weyl 中心随特征变化的机制
%

\subsection{量子平面的定义}
\label{b2:subsec:2003:01}

\begin{definition}\label{b2:def:quantum-plane}
设 \(k\) 为域，\(q\in k^\times\)。将自由结合代数 \(k\langle x,y\rangle\) 对关系 \(yx=qxy\) 生成的双边理想取商，所得 \(k\) 代数称为参数 \(q\) 的\term[liangzi-pingmian]{量子平面}[quantum plane]，记为
\[
k_q[x,y]=k\langle x,y\rangle/(yx-qxy).
\]
这里不要求 \(x,y\) 可逆。将它们也取逆得到的是另一种代数，不能与当前定义混同。
\end{definition}
\symidx{\(k_q[x,y]\)：参数 \(q\ne0\) 的量子平面}

\begin{theorem}\label{b2:thm:quantum-plane-normal-basis}
设 \(k\) 为域，\(q\in k^\times\)，\(k_q[x,y]=k\langle x,y\rangle/(yx-qxy)\) 为量子平面，则全部 \(x^iy^j\)、\(i,j\ge0\)，构成它的一组 \(k\) 基，且对所有非负整数 \(i,j,r,s\)，有
\begin{equation}\label{b2:eq:quantum-plane-product}
(x^iy^j)(x^ry^s)=q^{jr}x^{i+r}y^{j+s}.
\end{equation}
\end{theorem}
\begin{proof}
用 \(yx=qxy\) 移动所有错序字母，得到正规字的生成性。在乘积 \(x^iy^jx^ry^s\) 中，中间的每个 \(y\) 都要越过 \(r\) 个 \(x\)，总共交换 \(jr\) 次，每次产生一个因子 \(q\)。因而
\[
y^jx^r=q^{jr}x^ry^j,
\]
这给出了\eqref{b2:eq:quantum-plane-product}，但尚不能保证正规字独立。按照这个乘法取以 \(e_{ij}\)、\(i,j\ge0\) 为基的向量空间，令 \(e_{ij}\) 对应预期的正规字 \(x^iy^j\)，定义
\[
e_{ij}e_{rs}=q^{jr}e_{i+r,j+s}.
\]
乘三个基向量 \(e_{ij},e_{rs},e_{uv}\) 时，两种结合方式的系数指数分别为 \(jr+(j+s)u\) 与 \(su+j(r+u)\)，它们相同，所以乘法结合，单位为 \(e_{00}\)。元素 \(e_{10},e_{01}\) 满足量子平面的关系，故由\cref{b2:prop:algebra-relations-universal}得到到该代数的同态。正规字 \(x^iy^j\) 恰映到 \(e_{ij}\)，其独立性由后者为基得到。
\end{proof}

关系本身齐次，所以取 \(\deg x=\deg y=1\)，量子平面就是分次代数，第 \(n\) 次分量维数为 \(n+1\)。它与二元多项式环具有相同的分次维数，但乘法不同；例如 \(q\ne1\) 时，\(yx-xy=(q-1)xy\ne0\)。对这个分次所给的字长滤过取伴随分次，仍得到量子平面自身。Weyl 关系 \(\mathrm{d}\cdot x-x\cdot\mathrm{d}=1\) 则含低次修正，常数在最高次层消失，所以它的伴随分次是交换多项式环。量子关系中的 \(q\) 保留在同一次数内，取伴随分次不会消去这个扭曲。

\subsection{扭曲多项式环的实例}
\label{b2:subsec:2003:02}

令 \(\sigma:k[x]\to k[x]\)、\(\sigma\bigl(f(x)\bigr)=f(qx)\)。因为 \(q\ne0\)，它是自同构，逆为代入 \(q^{-1}x\)。由关系逐项计算，得到
\[
y\cdot f(x)=\sigma\bigl(f(x)\bigr)y,
\qquad y^r\cdot f(x)=\sigma^r\bigl(f(x)\bigr)y^r.
\]
所以每个元素可以唯一写为 \(\sum_j f_j(x)y^j\)，而乘法为
\[
\bigl(f(x)y^i\bigr)\bigl(g(x)y^j\bigr)
=f(x)\sigma^i\bigl(g(x)\bigr)y^{i+j}.
\]
这称为一个\term[niuqu-duoxiangshi-huan]{扭曲多项式环}[skew polynomial ring]，记为 \(k[x][y;\sigma]\)。在\secref{b2:sec:2005}中，会把它推广为允许导子修正的 Ore 扩张。

\begin{proposition}\label{b2:prop:quantum-plane-domain}
设 \(k\) 为域，\(q\in k^\times\)，\(k_q[x,y]=k\langle x,y\rangle/(yx-qxy)\) 为量子平面，则它没有零因子。若以正规式中 \(y\) 的最高指数为次数，任意两个非零元素乘积的次数等于两因子次数之和。
\end{proposition}
\begin{proof}
设两因子的最高项为 \(f_m(x)y^m\)、\(g_n(x)y^n\)。乘积的 \(y^{m+n}\) 系数为 \(f_m(x)\sigma^m\bigl(g_n(x)\bigr)\)。\(\sigma\) 是自同构，且 \(k[x]\) 是整环，所以这个系数非零。其他项的 \(y\) 次数更低，不能消去它，得到结论。
\end{proof}

例如 \(y(x^2+x+1)=(q^2x^2+qx+1)y\)。这里移动系数改变的是整段多项式的输入；不能只给整个多项式乘上同一个 \(q\)。若 \(q=-1\) 且特征不为二，偶次幂与 \(y\) 交换，奇次幂则改变符号。

\subsection{中心与参数的关系}
\label{b2:subsec:2003:03}

\begin{proposition}\label{b2:prop:quantum-plane-center}
设 \(k\) 为域，\(q\in k^\times\)，\(k_q[x,y]=k\langle x,y\rangle/(yx-qxy)\) 为量子平面。若 \(q\) 在 \(k^\times\) 中具有无限阶，则 \(Z\bigl(k_q[x,y]\bigr)=k\)。若 \(q\) 的阶为有限整数 \(\ell\ge1\)，则
\[
Z\bigl(k_q[x,y]\bigr)=k[x^\ell,y^\ell].
\]
后一情形中，量子平面作为中心上的模自由，基为 \(x^iy^j\)、\(0\le i,j<\ell\)，秩为 \(\ell^2\)。
\end{proposition}
\begin{proof}
由\eqref{b2:eq:quantum-plane-product}，正规式 \(P=\sum c_{ij}x^iy^j\) 与 \(x\) 交换恰好要求 \((q^j-1)c_{ij}=0\)，与 \(y\) 交换恰好要求 \((q^i-1)c_{ij}=0\)。正规单项式基使这些条件逐项成立而不能相互抵消。无限阶时只剩 \(i=j=0\)；有限阶时，\(q^r=1\) 恰好等价于 \(\ell\mid r\)，故要求 \(\ell\mid i,j\)，得到中心公式。

对一般指数写 \(i=a\ell+r,j=b\ell+s\)，其中 \(0\le r,s<\ell\)，再把中心因子提出，得到所列有限生成组。若它们有中心系数关系，展开这些系数后便得到不同的正规单项式之间的 \(k\) 线性关系，所有系数必须为零。因此它们确为中心上的基。
\end{proof}

这里的 \(\ell\) 必须是 \(q\) 的确切阶。仅有 \(q^\ell=1\) 可以保证 \(x^\ell,y^\ell\) 属于中心，却不能保证它们生成整个中心；例如 \(q=1\) 时虽然 \(q^2=1\)，中心仍为 \(k[x,y]\)，不是 \(k[x^2,y^2]\)。特别地，\(q=1\) 对应 \(\ell=1\)；\(q=-1\) 且特征不为二时，确切阶为二，中心为 \(k[x^2,y^2]\)，代数在中心上的秩为四。在特征二中，\(-1=1\)，同一写法则恢复交换多项式环。

\ProseParagraphBreak
量子平面中心的上述变化由 \(q\) 的乘法阶决定。\cref{b2:thm:weyl-center-simplicity}中 Weyl 代数的中心变化则由域特征决定：特征零时中心只有 \(k\)，特征 \(p>0\) 时中心为 \(k[x^p,\mathrm{d}^p]\)，在中心上的秩为 \(p^2\)。例如在特征零中取 \(q=-1\)，量子平面已在中心上有限，而 Weyl 代数仍有标量中心。两个有限秩公式都来自将正规单项式的指数按中心幂的指数取余，但使这些幂成为中心元的条件不同。

% !TeX root = ../../Book2.tex
% 纲要标题：### 20.4 Lie 代数与 PBW 型结构
% 纲要位置：计划纲要.md，第 1716 行。

% 纲要要点（供目录核对）：
%
% * 由交换的伴随分次控制非交换代数；含常数的一次滤过项、重排与独立性分开证明；Rees 的 Z² 来自齐次化最高次数，习题推广加权参数
% * Lie 代数的交替括号、Jacobi 恒等式、交换子与矩阵例子；在特征二保留交替性，区别 Lie 理想与结合代数双边理想
% * Lie 代数的表示及伴随表示；括号关系转为包络代数上的幺性左模，Heisenberg 的全参数中心商与正规基在原命题内证明
% * 与第五卷包络代数的 PBW 定理的联系；一般结论完整证明仍归第五卷第15章，j 的单射性只由 PBW 后取得，具体模型不借一般结论
% * 低次修正的 Jacobi 约束与 PBW 相容性；zyx 歧义与 Jacobi 差产生同一个隐藏关系，正文命题完整辨认实际二生成元商及正规基，原题改为加权 Ore 实例
%

\section{Lie 代数与 PBW 型结构}
\label{b2:sec:2004}

Weyl 代数的交换关系含常数修正，量子平面的扭曲关系则是齐次的。前两节分别借具体模型证明了有序单项式独立；若只给定一组带低次修正的关系，最高次交换关系只能给出从多项式环到伴随分次的满射，额外关系可能使它并非单射。结合代数的交换子满足 Jacobi 恒等式，Lie 代数将这一约束抽取为公理，泛包络代数再把括号写回结合乘法。有序单项式能否组成基，仍须由伴随分次检验。以下 \(k\) 为任意域。


\subsection{伴随分次与非交换结构}
\label{b2:subsec:2004:01}

\begin{proposition}\label{b2:prop:pbw-filtered-criterion}
设 \(k\) 为域，\(A\) 为由 \(x_1,\ldots,x_n\) 生成的含幺结合 \(k\) 代数，\(n\geq1\)。以这些生成元的字长不超过 \(m\) 的乘积张成 \(F_mA\)，其中 \(F_0A=k\cdot1\)、\(F_1A=k\cdot1+kx_1+\cdots+kx_n\)。若全部交换子 \([x_i,x_j]=x_ix_j-x_jx_i\) 均属于 \(F_1A\)，则有满同态
\[
k[X_1,\ldots,X_n]\longrightarrow\operatorname{gr}_F A,
\qquad X_i\longmapsto x_i+F_0A\in F_1A/F_0A.
\]
它是同构，当且仅当全部有序单项式 \(x_1^{a_1}\cdots x_n^{a_n}\) 是 \(A\) 的一组基。
\end{proposition}
\begin{proof}
交换子属于 \(F_1A\)，所以一次符号交换；字长滤过又使这些符号生成整个伴随分次，得到满同态。关系 \(x_jx_i=x_ix_j+[x_j,x_i]\) 的修正项是常数与生成元的线性组合。交换相邻错序字母时，二次项的字长不变而错序数减少；修正项的字长则严格减少。因此按字长及错序数的字典序归纳，可以把任意长度不超过 \(m\) 的字整理为全次数不超过 \(m\) 的有序单项式的线性组合。

若这些单项式已经是一组基，就连各个滤过项的基也被上述整理确定。相应次数的商片以相同次数的单项式符号为基，故满同态逐次单射。

反过来，若伴随分次同构于多项式环，任意非零的有限有序单项式关系中，取最高全次数的部分，其符号就是交换多项式单项式之间的非平凡关系，矛盾。因此有序单项式线性无关；生成性已证，得到基。
\end{proof}

满足这种有序基性质时，常称所讨论的构造具有 PBW 型性质。这里所指的是一个已经说明的基与滤过结论，而不是把任意非交换关系都统称为 PBW 定理。

\ProseParagraphBreak
Weyl 代数给出完全可算的形变。它的 Rees 代数由 \(X=xz,D=\mathrm{d}\cdot z,Z=z\) 生成。两个原生成元各乘了一次 \(z\)，所以
\[
DX-XD=(\mathrm{d}\cdot x-x\cdot\mathrm{d})z^2=z^2=Z^2.
\]
原关系的常数项必须补上两个 \(Z\)，才能与左侧同为二次。\(z\) 的中心性又给出
\[
ZX=XZ,\qquad ZD=DZ.
\]
这些关系给出整个代数。先将 \(Z\) 移到右侧，再用 \(DX=XD+Z^2\) 移动错序的 \(D,X\)；修正项使 \(X,D\) 的总个数减少二，因此按这个数和错序数归纳，可以把任何字写成 \(X^iD^jZ^r\) 的线性组合。它们在 Rees 代数中的像为 \(x^i\mathrm{d}^jz^{i+j+r}\)，由\cref{b2:thm:weyl-normal-basis}及不同 \(z\) 次数的独立性，它们线性无关。因此该关系代数与 Rees 代数同构。

\ProseParagraphBreak
它作为 \(k[Z]\) 模以 \(X^iD^j\) 为基。在 \(Z=0\) 处得到交换平面，在 \(Z=1\) 处得到 Weyl 代数；这是\cref{b2:thm:rees-fibres}的一个具体实现。

\subsection{Lie 代数与基本例子}
\label{b2:subsec:2004:lie-basics}

在 Weyl 代数中，交换子 \(\mathrm{d}\cdot x-x\cdot\mathrm{d}\) 比乘积 \(\mathrm{d}\cdot x\) 的次数低。现在暂时不保留两个乘积，只把它们的差作为一项双线性运算。结合律会使这种运算满足额外的恒等式；这些恒等式正是下面定义的依据。

\begin{definition}\label{b2:def:lie-algebra}
设 \(k\) 为域，\(\mathfrak g\) 为 \(k\) 向量空间，并给定双线性映射
\[
[\ ,\ ]:\mathfrak g\times\mathfrak g\longrightarrow\mathfrak g,
\]
若它满足交替性 \([u,u]=0\) 对所有 \(u\in\mathfrak g\) 成立，以及对任意 \(u,v,w\in\mathfrak g\) 成立的\term[Jacobi-hengdengshi]{Jacobi 恒等式}[Jacobi identity]：\footnote{雅可比（Carl Gustav Jacob Jacobi，1804-1851），德国数学家，研究椭圆函数、行列式和力学。}
\begin{equation}\label{b2:eq:lie-jacobi}
\bigl[u,[v,w]\bigr]+\bigl[v,[w,u]\bigr]+\bigl[w,[u,v]\bigr]=0.
\end{equation}
则称 \((\mathfrak g,[\ ,\ ])\) 为 \(k\) 上的\term[lie-daishu]{Lie 代数}[Lie algebra]，称 \([\ ,\ ]\) 为 Lie 括号。这里不要求结合律，也不指定单位元。
\end{definition}
\symidx{\([u,v]\)：Lie 代数的括号}

展开 \([u+v,u+v]=0\)，得 \([u,v]=-[v,u]\)。反过来，在特征不等于二时，反对称性推出 \([u,u]=0\)；特征二时这一推论失效，所以定义中保留交替性。这与\cref{b2:prop:alternating-implies-skew,b2:def:symmetric-alternating-form}中交替映射和双线性型的区别相同。例如一维空间 \(ke\) 上的双线性运算 \([e,e]=e\) 在特征二时反对称，却不交替。将任意向量空间上的括号恒取为零，则满足全部公理，称为\term[jiaohuan-lie-daishu]{交换 Lie 代数}[abelian Lie algebra]。

\begin{proposition}\label{b2:prop:commutator-lie-algebra}
设 \(k\) 为域，\(A\) 为结合 \(k\) 代数。在其底向量空间上定义 \([a,b]=ab-ba\) 对任意 \(a,b\in A\) 成立，则所得结构为 \(k\) 上的 Lie 代数，记为 \(A^-\)。
\end{proposition}
\begin{proof}
乘法的双线性给出括号的双线性，且 \([a,a]=0\)。对于 Jacobi 恒等式，逐项展开为
\[
\begin{aligned}
\bigl[a,[b,c]\bigr]&=abc-acb-bca+cba,\\
\bigl[b,[c,a]\bigr]&=bca-bac-cab+acb,\\
\bigl[c,[a,b]\bigr]&=cab-cba-abc+bac.
\end{aligned}
\]
结合律使这里的三重乘积无需另加括号；相加后每个乘积的两个系数互为相反数，所得和为零。这个消去在任意特征下成立。
\end{proof}
\symidx{\(A^-\)：结合代数的交换子 Lie 代数}

特别地，\(\mathfrak{gl}(V)=\operatorname{End}_k(V)^-\)，有限维且选定基时也记为 \(\mathfrak{gl}_n(k)=M_n(k)^-\)。其括号由算子复合或矩阵乘法之差给出，与群中的乘法交换子 \(aba^{-1}b^{-1}\) 是不同运算。
\symidx{\(\mathfrak{gl}(V)\)：向量空间的一般线性 Lie 代数}

\ProseParagraphBreak
设 \(\mathfrak g\) 为 \(k\) 上的 Lie 代数。若线性子空间 \(\mathfrak h\subseteq\mathfrak g\) 对括号封闭，则称 \(\mathfrak h\) 为\term[lie-zidaishu]{Lie 子代数}[Lie subalgebra]；公理可从 \(\mathfrak g\) 直接限制下来。对 \(k\) 上的 Lie 代数 \(\mathfrak g,\mathfrak h\)，若 \(k\) 线性映射 \(f:\mathfrak g\to\mathfrak h\) 保持括号，即
\[
f\bigl([u,v]\bigr)=\bigl[f(u),f(v)\bigr].
\]
则称 \(f\) 为\term[lie-daishu-tongtai]{Lie 代数同态}[Lie algebra homomorphism]。

\ProseParagraphBreak
若线性子空间 \(I\subseteq\mathfrak g\) 满足 \([u,a]\in I\) 对所有 \(u\in\mathfrak g\)、\(a\in I\) 成立，则称 \(I\) 为 \(\mathfrak g\) 的\term[lie-lixiang]{Lie 理想}[ideal of a Lie algebra]。交替性使 \([a,u]\) 也属于 \(I\)。这里的封闭条件只涉及括号，不能换成结合乘法的左右封闭性。例如在 \(M_2(k)^-\) 中，\(kI_2\) 是 Lie 理想，因为与任意矩阵的括号为零；它却不是 \(M_2(k)\) 的双边理想，因为 \(E_{11}I_2=E_{11}\notin kI_2\)。若 \(I\) 为 Lie 理想，则
\[
[u+I,v+I]=[u,v]+I
\]
在 \(\mathfrak g/I\) 上良定义：替换两个代表元所增加的 \([a,v]+[u,b]+[a,b]\) 均属于 \(I\)。双线性、交替性和 Jacobi 恒等式由商映射传递，因而得到商 Lie 代数。特别地，Lie 同态的核是这样的理想，其像是 Lie 子代数。

\ProseParagraphBreak
设 \(\mathfrak g\) 为 \(k\) 上的 Lie 代数。若 \(u\in\mathfrak g\) 与 \(\mathfrak g\) 的全部元素括号为零，则称 \(u\) 为中心元素。所有中心元素组成 \(\mathfrak g\) 的\term[lie-daishu-zhongxin]{中心}[center of a Lie algebra]，记为
\[
Z(\mathfrak g)=\bigl\{u\in\mathfrak g\mid [u,v]=0\;\text{对所有}\;v\in\mathfrak g\bigr\}.
\]
它是一个 Lie 理想，因为对任意 \(u\in Z(\mathfrak g)\)、\(v\in\mathfrak g\)，都有 \([v,u]=0\in Z(\mathfrak g)\)。
\symidx{\(Z(\mathfrak g)\)：Lie 代数的中心}

\ProseParagraphBreak
最小的非交换例子已能在二阶矩阵中看到。令 \(h=E_{11}\)、\(e=E_{12}\)，则
\[
[h,e]=e,\qquad [h,h]=[e,e]=0.
\]
因此 \(kh\oplus ke\) 对括号封闭，且
\[
[ah+be,ch+de]=(ad-bc)e.
\]
它是二维非交换 Lie 代数。在这里，封闭性和\cref{b2:prop:commutator-lie-algebra}已经保证 Jacobi 恒等式，不需要对每一组三元组重新检验。

\begin{example}\label{b2:exm:heisenberg-lie-algebra}
设 \(k\) 为域。\(M_3(k)^-\) 中三阶严格上三角矩阵组成的 Lie 代数记为 \(\mathfrak h\)，称为三维\term[heisenberg-lie-daishu]{Heisenberg Lie 代数}[Heisenberg Lie algebra]。\footnote{海森堡（Werner Karl Heisenberg，1901-1976），德国物理学家，创立矩阵力学，提出不确定性原理。}取
\[
x=E_{23},\qquad d=E_{12},\qquad z=E_{13},
\]
则 \(x,d,z\) 是它的一组基，括号满足
\begin{equation}\label{b2:eq:heisenberg-brackets}
[d,x]=z,\qquad [z,x]=[z,d]=0.
\end{equation}
确实，\(E_{ij}E_{rs}=\delta_{jr}E_{is}\) 给出 \(dx=z\)、\(xd=0\)，而 \(z\) 与 \(x,d\) 的两个方向乘积均为零。因而这些矩阵对括号封闭，\(z\) 是一个非零的中心元素。
\end{example}

上述交换子与矩阵例子也见 Etingof 等人的定义和例表\cite[Definition 2.9.1; Example 2.9.2, pp. 22--23]{EtingofEtAl2011RepresentationTheory}。这里先用矩阵检验抽象公理；稍后将通过同一个 Heisenberg 例子回到 Weyl 代数。

\subsection{Lie 代数的表示}
\label{b2:subsec:2004:lie-representations}

矩阵例子说明，Lie 代数可以通过线性算子的交换子实现。对于一个给定的抽象 Lie 代数，不必要求这种实现单射；它只是指定各元素怎样作用在某个向量空间上。

\begin{definition}\label{b2:def:lie-representation}
设 \(k\) 为域，\(\mathfrak g\) 为 \(k\) 上的 Lie 代数，\(V\) 为 \(k\) 向量空间。若 \(k\) 线性映射
\[
\rho:\mathfrak g\longrightarrow\operatorname{End}_k(V)
\]
对任意 \(u,v\in\mathfrak g\) 满足
\begin{equation}\label{b2:eq:lie-representation}
\rho\bigl([u,v]\bigr)=\rho(u)\rho(v)-\rho(v)\rho(u).
\end{equation}
则称 \(\rho\) 为 \(\mathfrak g\) 在 \(V\) 上的\term[lie-daishu-biaoshi]{表示}[representation of a Lie algebra]；换言之，\(\rho:\mathfrak g\to\mathfrak{gl}(V)\) 是 Lie 代数同态。记 \(u\cdot w=\rho(u)(w)\)，其条件就是
\[
[u,v]\cdot w=u\cdot(v\cdot w)-v\cdot(u\cdot w).
\]
带有这种作用的 \(V\) 也称为 \(\mathfrak g\) 模。
\termidx[lie-mo]{Lie 代数模}
\foreignidx{Lie algebra module}
\end{definition}

若对所有 \(u\) 取 \(\rho(u)=0\)，得到平凡表示。若 \(\mathfrak g\) 是矩阵 Lie 子代数，其矩阵对列向量的自然作用就是一个表示。若表示映射 \(\rho\) 单射，则称它为\term[zhongshi-biaoshi]{忠实表示}[faithful representation]。两个表示间的同态是满足 \(F(u\cdot v)=u\cdot F(v)\) 的线性映射 \(F:V\to W\)；它要求作用相容，不要求 \(F\) 可逆。

\begin{proposition}\label{b2:prop:adjoint-lie-representation}
设 \(k\) 为域，\(\mathfrak g\) 为 \(k\) 上的 Lie 代数。令
\[
\operatorname{ad}(u)(v)=[u,v].
\]
则 \(\operatorname{ad}:\mathfrak g\to\mathfrak{gl}(\mathfrak g)\) 是一个表示，称为\term[bansui-biaoshi]{伴随表示}[adjoint representation]。它的核是
\[
Z(\mathfrak g)=\bigl\{u\in\mathfrak g\mid [u,v]=0\;\text{对所有}\;v\in\mathfrak g\bigr\},
\]
即 \(\mathfrak g\) 的中心。
\end{proposition}
\begin{proof}
括号的双线性使 \(\operatorname{ad}(u)\) 为线性算子，并使 \(u\mapsto\operatorname{ad}(u)\) 线性。把\eqref{b2:eq:lie-jacobi}中的后两项用反对称性改写，得到
\[
\bigl[[u,v],w\bigr]
=\bigl[u,[v,w]\bigr]-\bigl[v,[u,w]\bigr].
\]
因此，对任意 \(w\)，
\[
\operatorname{ad}\bigl([u,v]\bigr)(w)
=\bigl(\operatorname{ad}(u)\operatorname{ad}(v)
-\operatorname{ad}(v)\operatorname{ad}(u)\bigr)(w),
\]
正是\eqref{b2:eq:lie-representation}。最后，\(\operatorname{ad}(u)=0\) 等价于 \([u,v]=0\) 对全部 \(v\) 成立，故其核恰为所写的中心。
\end{proof}
\symidx{\(\operatorname{ad}\)：Lie 代数的伴随表示}

这说明 Jacobi 恒等式控制着一个实际作用：先取括号 \([u,v]\) 再作用，与两个作用算子的交换子相同。它不保证伴随表示忠实。例如\cref{b2:exm:heisenberg-lie-algebra}中的 \(z\ne0\)，却有 \(\operatorname{ad}(z)=0\)。表示的定义与伴随表示也可对照 Etingof 等人\cite[Definition 2.9.7; Example 2.9.8, p. 24]{EtingofEtAl2011RepresentationTheory}。

\subsection{包络代数与 PBW 定理}
\label{b2:subsec:2004:02}

表示条件\eqref{b2:eq:lie-representation}混合了 Lie 括号与算子复合。若把每个 \(u\in\mathfrak g\) 当成一个结合代数的生成元，再加入恰好保证该条件的关系，就能用已经建立的模论讨论 Lie 表示。

\begin{definition}\label{b2:def:universal-enveloping-algebra}
设 \(k\) 为域，\(\mathfrak g\) 为 \(k\) 上的 Lie 代数，\(T_k(\mathfrak g)\) 为其底向量空间的张量代数。令 \(J\) 为全部元素
\[
u\otimes v-v\otimes u-[u,v],\qquad u,v\in\mathfrak g,
\]
生成的双边理想，其中最后一项属于一次张量分量。定义商代数
\[
U(\mathfrak g)=T_k(\mathfrak g)/J
\]
并称它为 \(\mathfrak g\) 的\term[fan-baoluo-daishu]{泛包络代数}[universal enveloping algebra]。以 \(j:\mathfrak g\to U(\mathfrak g)\) 表示一次分量包含与商映射的复合。
\end{definition}
\symidx{\(U(\mathfrak g)\)：Lie 代数的泛包络代数}

关系 \(u\otimes v-v\otimes u-[u,v]\) 的前两项属于张量次数二，最后一项属于张量次数一。因此张量字长仍给出商代数的滤过，原来的张量次数分次通常却不能下降到商；随后考察 \(\operatorname{gr}U(\mathfrak g)\)，正是为了分离这些关系的最高次部分。

\ProseParagraphBreak
定义只给出线性映射 \(j\)，尚未证明它单射。在证明单射以前，只能使用其像 \(j(\mathfrak g)\)，不能把 \(\mathfrak g\) 直接视为 \(U(\mathfrak g)\) 的线性子空间。商关系保证
\begin{equation}\label{b2:eq:enveloping-canonical-lie-map}
j\bigl([u,v]\bigr)=j(u)j(v)-j(v)j(u),
\end{equation}
所以 \(j:\mathfrak g\to U(\mathfrak g)^-\) 是 Lie 同态。这个构造与\cref{b2:def:enveloping-algebra}的 \(A^e=A\otimes A^{\mathrm{op}}\) 不同：\(A^e\) 用于将结合代数的双模视为左模，而 \(U(\mathfrak g)\) 用于将 Lie 代数的表示视为结合代数的模。

\begin{theorem}[泛包络代数的泛性质]\label{b2:thm:enveloping-universal}
设 \(k\) 为域，\(\mathfrak g\) 为 \(k\) 上的 Lie 代数，\(U(\mathfrak g)\) 为其泛包络代数，\(j:\mathfrak g\rightarrow U(\mathfrak g)\) 为自然 Lie 同态。对任意含幺结合 \(k\) 代数 \(B\) 和 Lie 同态 \(f:\mathfrak g\to B^-\)，其中 \(B^-\) 使用交换子括号，存在唯一保幺的 \(k\) 代数同态
\[
\widetilde f:U(\mathfrak g)\longrightarrow B,
\qquad \widetilde f\circ j=f.
\]
反过来，每个这样的代数同态与 \(j\) 复合，都给出一个 Lie 同态。
\end{theorem}
\begin{proof}
先将 \(f\) 视为线性映射。\cref{b2:thm:tensor-algebra-universal}给出唯一代数同态 \(\widehat f:T_k(\mathfrak g)\to B\)，在一次分量上等于 \(f\)。对关系理想的每个指定生成元，有
\[
\widehat f\bigl(u\otimes v-v\otimes u-[u,v]\bigr)
=f(u)f(v)-f(v)f(u)-f\bigl([u,v]\bigr)=0,
\]
最后的等式正是 \(f\) 保持括号的条件。\(\widehat f\) 是代数同态，所以也零化这些生成元的任意左右倍数之和，即整个 \(J\)。于是 \(\widetilde f(t+J)=\widehat f(t)\) 良定义，给出所需代数同态。\(U(\mathfrak g)\) 由 \(j(\mathfrak g)\) 及单位元生成，故其像一经指定，同态唯一。

反向若给定 \(\widetilde f\)，将它作用于\eqref{b2:eq:enveloping-canonical-lie-map}，便得 \(\widetilde f\circ j\) 保持括号。由刚证的唯一性，两个过程互逆。
\end{proof}

取零 Lie 同态 \(\mathfrak g\to k^-\)，得到代数同态 \(U(\mathfrak g)\to k\)，它将每个 \(j(u)\) 送到零，将单位元送到一。因此 \(U(\mathfrak g)\) 的单位元非零。

\begin{corollary}\label{b2:cor:lie-representations-enveloping-modules}
设 \(k\) 为域，\(\mathfrak g\) 为 \(k\) 上的 Lie 代数，\(V\) 为 \(k\) 向量空间，\(U(\mathfrak g)\) 为泛包络代数，则 \(\mathfrak g\) 在 \(V\) 上的表示与 \(V\) 上的幺性左 \(U(\mathfrak g)\) 模结构一一对应；对任意两个这样的向量空间，Lie 表示同态恰为相应的 \(U(\mathfrak g)\) 模同态。
\end{corollary}
\begin{proof}
将\cref{b2:thm:enveloping-universal}用于 \(B=\operatorname{End}_k(V)\)，Lie 表示 \(\rho\) 唯一延拓为代数同态 \(\widetilde\rho:U(\mathfrak g)\to\operatorname{End}_k(V)\)。规定 \(a\cdot v=\widetilde\rho(a)(v)\)，乘法与单位元的保持性恰给出幺性左模。反向把一个幺性左模的作用限制到 \(j(\mathfrak g)\)，由\eqref{b2:eq:enveloping-canonical-lie-map}得到 Lie 表示；唯一性说明两个过程互逆。

若线性映射 \(F:V\to W\) 与每个 \(u\in\mathfrak g\) 的作用相容，对乘积的长度归纳便知它与每个 \(j(u_1)\cdots j(u_r)\) 的作用相容；线性性再使它与这些乘积及单位元的任意线性组合相容，即为 \(U(\mathfrak g)\) 模同态。反向的限制立即给出 Lie 表示同态。
\end{proof}

表示中的括号相容条件被写成了 \(U(\mathfrak g)\) 的代数关系，而连续施加 Lie 元素的作用成为这个代数中的乘法。由于 \(U(\mathfrak g)\) 由 \(j(\mathfrak g)\) 和单位元生成，一个子空间对所有 Lie 元素的作用稳定，恰好等价于它是 \(U(\mathfrak g)\) 子模。于是前面的子模、商模和自同态环理论都可以用于 Lie 表示。

\ProseParagraphBreak
在交换 Lie 代数的情形，关系只有 \(uv-vu=0\)，故由\cref{b2:def:symmetric-algebra}，\(U(\mathfrak g)=\operatorname{Sym}_k(\mathfrak g)\)。若 \(\mathfrak g\) 有有限基 \(x_1,\ldots,x_n\)，\cref{b2:thm:symmetric-polynomial-algebra}将它识别为 \(k[x_1,\ldots,x_n]\)。因而交换 Lie 代数的表示就是一组两两交换的线性算子。

\begin{proposition}\label{b2:prop:heisenberg-enveloping-weyl}
设 \(k\) 为域，\(\mathfrak h\) 为以 \(x,d,z\) 为基、满足 \([d,x]=z\)、\([z,x]=[z,d]=0\) 的三维 Heisenberg Lie 代数，\(j:\mathfrak h\rightarrow U(\mathfrak h)\) 为自然映射。对任意 \(\lambda\in k\)，记
\[
B_\lambda=U(\mathfrak h)/\bigl(j(z)-\lambda1\bigr),
\]
其中括号表示所列元素生成的双边理想，\(\bar x,\bar d\) 分别为 \(j(x),j(d)\) 在 \(B_\lambda\) 中的像。\(B_\lambda\) 的 \(k\) 基为全部 \(\bar x^i\bar d^j\)、\(i,j\ge0\)。当 \(\lambda=0\) 时，\(B_0\cong k[X,D]\)，同构将 \(\bar x,\bar d\) 送到 \(X,D\)；当 \(\lambda\ne0\) 时，\(B_\lambda\cong A_1(k)\)，同构将 \(\bar x,\bar d\) 分别送到 Weyl 生成元 \(x,\lambda\cdot\mathrm{d}\)。对任意 \(k\) 向量空间 \(V\)，其上的幺性左 \(B_\lambda\) 模结构恰好对应中心元素 \(z\) 作用为 \(\lambda\operatorname{id}_V\) 的 \(\mathfrak h\) 表示。
\end{proposition}
\begin{proof}
商关系给出 \(\bar d\bar x-\bar x\bar d=\lambda1\)，并使 \(j(z)\) 的像为 \(\lambda1\)，所以 \(B_\lambda\) 由 \(\bar x,\bar d\) 及单位元生成。

若 \(\lambda=0\)，这两个生成元交换，故 \(X\mapsto\bar x,D\mapsto\bar d\) 给出代数同态 \(k[X,D]\to B_0\)。反向指定 Lie 同态 \(\mathfrak h\to k[X,D]^-\)，将 \(x,d,z\) 分别送到 \(X,D,0\)。因为目标的括号为零，这些像满足\eqref{b2:eq:heisenberg-brackets}，由\cref{b2:thm:enveloping-universal}得到 \(U(\mathfrak h)\to k[X,D]\)，再由 \(j(z)\) 映为零下降到 \(B_0\)。两个复合均固定生成元和单位元，因而互为逆映射。交换多项式环的基给出 \(B_0\) 的所列基。

若 \(\lambda\ne0\)，在 \(A_1(k)^-\) 中指定 \(x\mapsto x,d\mapsto\lambda\cdot\mathrm{d},z\mapsto\lambda1\)。等式 \([\lambda\cdot\mathrm{d},x]=\lambda1\) 和标量的中心性保证这是 Lie 同态；同样的泛性质和取商过程给出 \(B_\lambda\to A_1(k)\)。反向，\(\lambda^{-1}\bar d\) 与 \(\bar x\) 满足 Weyl 关系，故由\cref{b2:prop:algebra-relations-universal}得到 \(A_1(k)\to B_\lambda\)，将 \(x,\mathrm{d}\) 分别送到 \(\bar x,\lambda^{-1}\bar d\)。两个复合再次固定生成元和单位元，所以互为逆映射。所列单项式在前一个同构下成为 \(\lambda^jx^i\mathrm{d}^j\)，其中 \(\lambda^j\ne0\)，由\cref{b2:thm:weyl-normal-basis}得到基。

最后，\cref{b2:cor:lie-representations-enveloping-modules}将 \(\mathfrak h\) 表示视为 \(U(\mathfrak h)\) 模，而这种作用下降到 \(B_\lambda\) 当且仅当 \(j(z)-\lambda1\) 的作用为零，即 \(z\) 作用为 \(\lambda\operatorname{id}_V\)。
\end{proof}

取 \(\lambda=1\) 就得到原来的 Weyl 代数；这个辨认可与 Etingof 等人的 Heisenberg 例子对照\cite[Example 2.9.13, p. 24]{EtingofEtAl2011RepresentationTheory}。取 \(\lambda=0\) 则得到交换多项式环。同一个 Lie 代数的中心元素在表示中取不同的标量作用，所对应的结合代数商可以由非交换变为交换。

\ProseParagraphBreak
对有限维 \(\mathfrak g\) 选有序基 \(x_1,\ldots,x_n\)，令 \(F_mU(\mathfrak g)\) 为长度不超过 \(m\) 的 \(j(x_i)\) 乘积所生成的线性空间。\eqref{b2:eq:enveloping-canonical-lie-map}的右侧交换子属于一次滤过项，所以\cref{b2:prop:pbw-filtered-criterion}给出自然满射
\[
\operatorname{Sym}_k(\mathfrak g)\longrightarrow\operatorname{gr}U(\mathfrak g),
\qquad x_i\longmapsto j(x_i)+F_0U(\mathfrak g).
\]
这里还使用\cref{b2:thm:symmetric-polynomial-algebra}把给定基上的对称代数与多项式环识别。满射的单射性是下面具名结果的内容。

\begin{claim}[PBW]\label{b2:claim:pbw-enveloping-preview}
设 \(k\) 为域，\(\mathfrak g\) 为\(k\) 上的有限维 Lie 代数，\(x_1,\ldots,x_n\) 为其有序基，\(U(\mathfrak g)\) 为泛包络代数，\(j:\mathfrak g\rightarrow U(\mathfrak g)\) 为自然映射。按 \(j(x_i)\) 的乘积长度滤过 \(U(\mathfrak g)\)，则自然映射 \(\operatorname{Sym}_k(\mathfrak g)\rightarrow\operatorname{gr}U(\mathfrak g)\) 为同构；等价地，全部 \(j(x_1)^{a_1}\cdots j(x_n)^{a_n}\)、\(a_i\geq0\)，构成 \(U(\mathfrak g)\) 的 \(k\) 基。特别地，\(j\) 单射。
\end{claim}

这称为\term[PBW-dingli]{Poincaré--Birkhoff--Witt 定理}[Poincaré--Birkhoff--Witt theorem]。\footnote{庞加莱（Henri Poincaré，1854-1912），法国数学家，研究拓扑、动力系统与微分方程。伯克霍夫（Garrett Birkhoff，1911-1996），美国数学家，研究格论与普遍代数。Witt 已在本卷前文介绍。}完整证明见第五卷第15章。前面的 Weyl 与量子平面正规基，以及上述 Heisenberg 商的基，已经由具体模型或互逆同态证明。较低次的括号关系必须满足相容条件，Jacobi 恒等式正是其中的关键。

\subsection{低次修正与 PBW 相容性}
\label{b2:subsec:2004:03}

考虑三个生成元 \(x,y,z\)，随意指定
\[
yx=xy+x,\qquad zy=yz+y,\qquad zx=xz.
\]
这些关系的最高次部分都是交换关系，但不保证基仍是 \(x^ay^bz^c\)。

\begin{proposition}\label{b2:exm:incompatible-pbw-relations}
设 \(k\) 为域，并令
\[
A=k\langle x,y,z\rangle/(yx-xy-x,\ zy-yz-y,\ zx-xz).
\]
在 \(A\) 中必有 \(x=0\)，且将 \(y,z\) 的同名商类对应起来给出代数同构
\[
A\cong k\langle y,z\rangle/(zy-yz-y).
\]
全部 \(y^iz^j\)、\(i,j\ge0\)，构成 \(A\) 的 \(k\) 基。因此预期的三变量 PBW 基虽然失败，余下两个生成元仍有唯一的正规式。
\end{proposition}
\begin{proof}
对同一个字 \(zyx\) 先整理左边的 \(zy\)，得到
\[
(zy)x=(yz+y)x=yxz+yx=xyz+xz+xy+x.
\]
先整理右边的 \(yx\)，则得到
\[
z(yx)=z(xy+x)=xzy+xz=xyz+xy+xz.
\]
结合律使两者相等，故 \(x=0\)。这里出现的是一次额外关系，单看最高次的三条交换关系无法预见。

记 \(B=k\langle y,z\rangle/(zy-yz-y)\)。将 \(x,y,z\) 分别送到 \(0,y,z\)，\(A\) 的三条定义关系均成立，得到同态 \(A\to B\)。反向，\(A\) 中的 \(y,z\) 满足 \(B\) 的定义关系，给出同态 \(B\to A\)。它们的复合固定 \(y,z\) 和单位元；在 \(A\) 中又有 \(x=0\)，所以两个复合处处为恒等，得到所述同构。

在 \(B\) 中，用 \(zy=yz+y\) 整理错序字母，二次项的错序数下降，一次项的字长下降，故按字长及错序数归纳，\(y^iz^j\) 张成 \(B\)。为了检验独立性，取交换多项式空间 \(k[Y,Z]\)，以 \(M_Y\) 表示乘以 \(Y\) 的算子。用 \(Y\partial_Y\) 产生关系中的 \(y\)，再加上与 \(M_Y\) 交换的乘 \(Z\) 算子，以便作用于 \(1\) 时仍能记录 \(z\) 的指数。具体令
\[
H(f)=Y\frac{\partial f}{\partial Y}+Zf.
\]
形式求导的 Leibniz 公式在任意特征下都给出
\[
(HM_Y-M_YH)(f)=Yf=M_Y(f).
\]
因此 \(y\mapsto M_Y,z\mapsto H\) 满足 \(B\) 的关系，得到代数同态 \(B\to\operatorname{End}_k(k[Y,Z])\)。由于 \(H(Z^j)=Z^{j+1}\)，有 \(H^j(1)=Z^j\)，从而
\[
M_Y^iH^j(1)=Y^iZ^j.
\]
若正规单项式间存在有限线性关系，将相应算子作用于 \(1\)，便得到不同交换单项式 \(Y^iZ^j\) 间的同一线性关系，所有系数只能为零。这证明了 \(B\) 以及 \(A\) 的所列基。
\end{proof}

这种失败可以用括号解释：指定的规则给出 \([y,x]=x\)、\([z,y]=y\)、\([z,x]=0\)。在以 \(x,y,z\) 为基的三维空间上，它们使
\[
\bigl[z,[y,x]\bigr]+\bigl[y,[x,z]\bigr]+\bigl[x,[z,y]\bigr]=-x\ne0,
\]
所以不满足 Jacobi 恒等式。这个 Jacobi 差为 \(-x\)，正是前面两种整理结果之差 \(x\) 的相反数，二者暴露的是同一个额外关系。\cref{b2:prop:commutator-lie-algebra}已证明，结合代数中的加法交换子必满足该恒等式，因此关系只能通过使 \(x\) 消失来相容。\secref{b2:sec:2006}将把这种同一字的两种整理方式作为重叠歧义处理。

% !TeX root = ../../Book2.tex
% 纲要标题：### 20.5* Ore 扩张与微分算子环
% 纲要位置：计划纲要.md，第 1724 行。

\OptionalSection{Ore 扩张与微分算子环}{b2:sec:2005}

Weyl 关系在移动系数时增加一个导数项，量子平面则对系数施加自同构。Ore 扩张把这两种操作合在一起，添加一个新的形式变量，并规定它与原环系数怎样相乘。第十九章的 Ore 局部化则把原环中指定的元素变成可逆元；这里的新变量并不因构造而可逆。

% 纲要要点（供目录核对）：
%
% * 环自同态 \(\sigma\) 与 \(\sigma\)-导子 \(\delta\)：\(\delta(ab)=\sigma(a)\delta(b)+\delta(a)b\)；结合律强迫扭曲 Leibniz 条件
% * Ore 扩张 \(R[t;\sigma,\delta]\) 的关系 \(ta=\sigma(a)t+\delta(a)\) 与正规表达；自由左模上的加法算子在一上读取系数，正文命题证明自同构时的右正规式及非单射边界
% * 微分域上的 \(K[\partial;\delta]\)：\(\partial a=a\partial+\delta(a)\)；区别形式生成元与其作用，Weyl 嵌入由正规式而非关系成立保证
% * 左模与线性微分方程；左理想、首一左除法和解同态就地证明，伴随矩阵的模坐标与取值转置写为正式命题；原题改为含系数导数的换基，第三卷附录 E 和第五卷附录 E 作为延伸
%

\subsection{环自同态与扭曲导子}
\label{b2:subsec:2005:01}

\begin{definition}\label{b2:def:sigma-derivation}
设 \(R\) 为含幺环，\(\sigma:R\to R\) 为保幺环自同态（不要求可逆），\(\delta:R\to R\) 为加法映射。若对任意 \(a,b\in R\)，有
\[
\delta(ab)=\sigma(a)\delta(b)+\delta(a)b,
\]
则称 \(\delta\) 为\term[sigma-daozi]{左 \(\sigma\) 导子}[left sigma-derivation]。这里的“左”指 \(\sigma\) 作用于乘积的左因子；\(\sigma=\operatorname{id}\) 时就是通常的导子。
\end{definition}

由 \(\delta(1)=\delta(1)+\delta(1)\) 得 \(\delta(1)=0\)。例如 \(k[x]\) 上的形式求导是恒等自同态的导子；对于任意自同态，零映射总是一个 \(\sigma\) 导子。

\ProseParagraphBreak
扭曲导子公式是结合律 \((ta)b=t(ab)\) 强迫的相容条件。若希望 \(ta=\sigma(a)t+\delta(a)\)，则先移动 \(a\) 再移动 \(b\) 会得到
\[
(ta)b=\sigma(a)\sigma(b)t+\sigma(a)\delta(b)+\delta(a)b.
\]
它须与 \(t(ab)\) 的展开相同，因而需要上述 \(\sigma\) 导子等式。不过检查这一处相容还不是整个代数存在性的证明，下面给出一个实际模型。

\subsection{Ore 扩张的关系与正规表达}
\label{b2:subsec:2005:02}

\begin{theorem}\label{b2:thm:ore-extension-construction}
设 \(R\) 为含幺环，\(\sigma:R\rightarrow R\) 为保幺环自同态，\(\delta:R\rightarrow R\) 为满足 \(\delta(ab)=\sigma(a)\delta(b)+\delta(a)b\) 的加法映射，则存在含幺环 \(R[t;\sigma,\delta]\)，含 \(R\) 为子环，且每个元素唯一写成有限和
\[
a_0+a_1t+\cdots+a_nt^n,\qquad a_i\in R,
\]
满足 \(ta=\sigma(a)t+\delta(a)\)。它作为左 \(R\) 模以 \(1,t,t^2,\ldots\) 为基。

对任意含幺环 \(B\)、保幺环同态 \(f:R\to B\) 及满足下式对所有 \(a\in R\) 成立的 \(u\in B\)，
\[
u\cdot f(a)=f\bigl(\sigma(a)\bigr)u+f\bigl(\delta(a)\bigr)
\]
存在唯一保幺环同态 \(R[t;\sigma,\delta]\to B\)，延伸 \(f\) 并把 \(t\) 送到 \(u\)。
\end{theorem}
\begin{proof}
为证明各个 \(t^i\) 左线性无关，取形式左自由模 \(P=\bigoplus_{n\ge0}Rz^n\)，并寻找一个算子 \(T\)，使 \(T^i(1)=z^i\)。这样，把候选关系作用于一，就能读出每个幂的系数。同时，\(T\) 必须实现题设的扭曲系数关系。因此在 \(P\) 的加法群上定义
\[
L_a\left(\sum_n b_nz^n\right)=\sum_n ab_nz^n,
\]
\[
T\left(\sum_n b_nz^n\right)
=\sum_n\bigl(\sigma(b_n)z^{n+1}+\delta(b_n)z^n\bigr).
\]
它们是加法群自同态。对 \(b\in R,p\in P\)，导子公式给出
\[
T(bp)=\sigma(b)T(p)+\delta(b)p,
\]
所以 \(T\) 通常不是左 \(R\) 线性的；非交换环中，\(L_a\) 也不必左 \(R\) 线性。这就是在 \(\operatorname{End}_{\mathbb Z}(P)\) 中构造模型，而不使用 \(\operatorname{End}_R(P)\) 的原因。直接在 \(bz^n\) 上计算得
\[
TL_a=L_{\sigma(a)}T+L_{\delta(a)}.
\]
令 \(S\) 为 \(\operatorname{End}_{\mathbb Z}(P)\) 中由全部 \(L_a\) 和 \(T\) 生成的子环。利用刚才的等式，可以在每个有限算子字中把所有 \(L_a\) 移到左侧；对 \(T^iL_a\) 按 \(i\) 归纳即可看到，每次展开仍为有限和。因此 \(S\) 中元素都具有 \(\sum_iL_{a_i}T^i\) 的形式。

由于 \(\sigma(1)=1\)、\(\delta(1)=0\)，有 \(T^i(1)=z^i\)。若 \(\sum_iL_{a_i}T^i=0\)，作用在一上便得 \(\sum_i a_iz^i=0\)，所有系数为零。这同时证明正规式唯一，以及 \(a\mapsto L_a\) 单射。令 \(R[t;\sigma,\delta]=S\)、\(t=T\)，结合律由算子复合继承。

对于泛性质，正规式唯一使公式 \(\sum_i a_it^i\mapsto\sum_i f(a_i)u^i\) 良定义。计算两个正规式的乘积时，逐次用 \(ta=\sigma(a)t+\delta(a)\) 移动系数；将每一步的 \(a,t\) 换成 \(f(a),u\)，题设等式保证同一步展开仍成立。因此该公式保持乘法，也保持加法与单位。任何延伸同态的值都由 \(R\) 及 \(t\) 确定，故唯一。
\end{proof}

这个环称为\term[Ore-kuozhang]{Ore 扩张}[Ore extension]。\(\delta=0\) 时简记为 \(R[t;\sigma]\)；\(\sigma=\operatorname{id}\) 时简记为 \(R[t;\delta]\)。上述定理给出的是系数在左侧的正规式；右侧的表达 \(\sum_i t^ib_i\) 是否存在、是否唯一，还须另证。\cref{b2:prop:ore-right-normal-form}将证明自同构情形的右正规式，并给出一般自同态情形的失败例子。
\symidx{\(R[t;\sigma,\delta]\)：关系 \(ta=\sigma(a)t+\delta(a)\) 的 Ore 扩张}

\begin{proposition}\label{b2:prop:ore-leading-degree}
设 \(R\) 为含幺环，\(\sigma:R\rightarrow R\) 为保幺环自同态，\(\delta\) 为 \(\sigma\) 导子，\(R[t;\sigma,\delta]\) 为相应 Ore 扩张。若非零元素 \(f,g\) 的最高项分别为 \(a_mt^m,b_nt^n\)，则乘积的 \(t^{m+n}\) 系数为 \(a_m\sigma^m(b_n)\)。因此，若 \(R\) 无零因子且 \(\sigma\) 单射，该 Ore 扩张也无零因子，非零乘积的次数相加。
\end{proposition}
\begin{proof}
由基本关系归纳得 \(t^m b=\sigma^m(b)t^m+\text{较低次项}\)。所以唯一可能贡献最高次项的是两个最高项的乘积，系数正如所述。单射性保证 \(\sigma^m(b_n)\ne0\)，再由 \(R\) 无零因子得到结论。
\end{proof}

不假设 \(\sigma\) 单射时，左正规式仍唯一，但环可以有零因子。例如 \(R=k[x]\)、\(\sigma(f)=f(0)\)、\(\delta=0\) 满足所有构造条件，却有 \(tx=0\)，而 \(t,x\) 都由正规式定理保证非零；另一方面，左正规单项式 \(xt\ne0\)。因此左、右乘法的行为不能从这条左正规式定理推成对称的结论。

\subsection{微分算子环与 Weyl 代数}
\label{b2:subsec:2005:03}

设 \(K\) 为带导子 \(\delta\) 的交换域。定义形式\term[weifen-suanzi-huan]{微分算子环}[differential operator ring]
\[
D=K[\partial;\delta],\qquad
\partial a=a\partial+\delta(a).
\]
元素为 \(\sum_i a_i\partial^i\)，乘法表示先作用右侧算子，再作用左侧算子。其在 \(K\) 上的自然作用为
\[
\left(\sum_i a_i\partial^i\right)(f)=\sum_i a_i\delta^i(f).
\]
这里 \(\partial\) 是 \(D\) 中的形式生成元，\(\delta\) 则是它在 \(K\) 上的作用算子。Ore 泛性质给出保幺同态 \(D\to\operatorname{End}_{\mathbb Z}(K)\)，保证这个作用与乘法相容，但不保证忠实。例如 \(\delta=0\) 时，形式变量 \(\partial\ne0\)，却在 \(K\) 上作用为零。

\ProseParagraphBreak
取 \(K=k(x)\)、\(\delta=\mathrm{d}/\mathrm{d}x\)，便有嵌入 \(A_1(k)\to K[\partial;\delta]\)，把 \(x,\mathrm{d}\) 送到 \(x,\partial\)。关系成立保证同态存在，Ore 正规式又保证 \(\sum_j f_j(x)\partial^j=0\) 时每个多项式系数都为零，所以它单射。这把多项式系数扩大为有理函数系数。正特征中仍应区别抽象形式算子与它在函数上的具体作用。

\ProseParagraphBreak
例如当 \(\delta=\mathrm{d}/\mathrm{d}x\) 时，
\[
\partial^2 a=a\partial^2+2\delta(a)\partial+\delta^2(a).
\]
这由连续两次移动系数得到。若把 \(a\) 直接与 \(\partial\) 交换，就会同时漏掉中间项和最后一项。

\subsection{左模与线性微分方程}
\label{b2:subsec:2005:04}

在 \(K\) 向量空间 \(M\) 上指定左 \(D\) 模结构，等价于指定一个加法算子 \(\nabla:M\to M\)，满足
\[
\nabla(am)=a\nabla(m)+\delta(a)m.
\]
正向取 \(\nabla(m)=\partial m\)，反向把标量作用及 \(\nabla\) 代入\cref{b2:thm:ore-extension-construction}的泛性质。配有这个算子的空间称为相对于 \(\delta\) 的\term[weifen-mo]{微分模}[differential module]。模同态就是与 \(\nabla\) 交换的 \(K\) 线性映射。\(\nabla\) 通常不是 \(K\) 线性的。

\ProseParagraphBreak
若要用循环生成元 \(e\) 记录方程 \(Le=0\)，左模作用还要求每个 \(QL\) 都消去 \(e\)。因此应模掉 \(L\) 生成的左理想 \(DL\)，从而构造左 \(D\) 模；不能无条件改用右理想 \(LD\)。

\begin{proposition}\label{b2:prop:differential-cyclic-module}
设 \(K\) 为带导子 \(\delta:K\rightarrow K\) 的交换域，\(D=K[\partial;\delta]\) 为满足 \(\partial a=a\partial+\delta(a)\) 的微分算子环，并让 \(D\) 按 \(\partial(y)=\delta(y)\) 作用于 \(K\)。取整数 \(m\geq1\) 及 \(a_0,\ldots,a_{m-1}\in K\)，令 \(L=\partial^m+a_{m-1}\partial^{m-1}+\cdots+a_0\)、\(M_L=D/DL\)，其中 \(DL\) 是左理想，则 \(M_L\) 有左 \(K\) 基 \(1,\partial,\ldots,\partial^{m-1}\) 的陪集，且有自然双射
\[
\operatorname{Hom}_D(M_L,K)\longrightarrow\bigl\{\,y\in K\mid L(y)=0\,\bigr\},
\qquad f\longmapsto f(1+DL).
\]
\end{proposition}
\begin{proof}
对任意算子 \(P\)，若最高项为 \(b_n\partial^n\)、\(n\ge m\)，就减去左倍数 \(b_n\partial^{n-m}L\)。由于 \(L\) 首一，这个倍数的最高项正是 \(b_n\partial^n\)，不必再求逆某个最高系数。每步严格降低次数，所以得到 \(P=QL+R\)，其中 \(R=0\) 或 \(\deg R<m\)。若有两个这样的余项，其差属于 \(DL\)；由\cref{b2:prop:ore-leading-degree}，非零 \(QL\) 的次数至少为 \(m\)，故余项之差只能为零。这证明所列商类为基。

左模同态由 \(1+DL\) 的像 \(y\) 决定，并必须满足 \(L(y)=0\)。反过来，若该方程成立，公式 \(P+DL\mapsto P(y)\) 与代表元无关，因为 \((QL)(y)=Q\bigl(L(y)\bigr)=0\)；它直接保持左 \(D\) 作用。两种构造互逆。
\end{proof}

\(M_L\) 记录的是方程对一个循环生成元施加的关系；解对应从 \(M_L\) 到函数模 \(K\) 的 \(D\) 模同态，并不是 \(M_L\) 的所有元素。解空间自然是常数域 \(C=\bigl\{\,a\in K\mid\delta(a)=0\,\bigr\}\) 上的向量空间：对 \(c\in C\)，有 \(\delta^i(cy)=c\delta^i(y)\)，故 \(L(cy)=cL(y)\)。一般 \(K\) 系数却会产生额外的导数项。例如方程 \(\partial y=0\) 的解包含一，但 \(\delta(a)\ne0\) 时 \(a\cdot1\) 不是解，所以不能一般地把解空间视为 \(K\) 子空间。\(C\) 确为域：导子公式保证对和、积封闭，并由 \(\delta(aa^{-1})=0\) 得到非零常数的逆仍为常数。

\ProseParagraphBreak
\(M_L\) 的 \(K\) 维数为 \(m\)，并不保证在指定的微分域中能找到 \(m\) 个独立解。例如取 \(K=\mathbb C(x)\)、\(\delta=\mathrm{d}/\mathrm{d}x\)、\(L=\partial-1\)。上述命题给出 \(\dim_KM_L=1\)，但 \(y'=y\) 没有非零有理函数解。若 \(y=p/q\ne0\)，其中 \(p,q\in\mathbb C[x]\) 均非零，则方程要求
\[
p'q-pq'=pq.
\]
左边为零或次数至多为 \(\deg p+\deg q-1\)，右边的次数却是 \(\deg p+\deg q\)，矛盾。因此这里 \(\operatorname{Hom}_D(M_L,K)=0\)。方程所对应的模已构造出来，解空间是否足够大还取决于允许解落在哪个微分域中。

\begin{proposition}\label{b2:prop:differential-companion-coordinates}
设 \(K\) 为带导子 \(\delta:K\to K\) 的交换域，\(D=K[\partial;\delta]\) 满足 \(\partial a=a\partial+\delta(a)\)，并按 \(\partial(y)=\delta(y)\) 左作用于 \(K\)。取整数 \(m\ge1\) 及 \(a_0,\ldots,a_{m-1}\in K\)，令
\[
L=\partial^m+\sum_{i=0}^{m-1}a_i\partial^i,\qquad
M_L=D/DL,\qquad e_i=\partial^i+DL\quad(0\le i<m).
\]
以 \(e_0,\ldots,e_{m-1}\) 为左 \(K\) 基，将元素写成坐标列 \(v=(v_0,\ldots,v_{m-1})^{\mathsf T}\)。令 \(A\in M_m(K)\) 的行、列按 \(0,\ldots,m-1\) 编号，条目为
\[
A_{ij}=
\begin{cases}
1,&i=j+1,\quad 0\le j<m-1,\\
-a_i,&j=m-1,\\
0,&\text{其余情形}.
\end{cases}
\]
则左乘 \(\partial\) 的坐标为 \(\nabla(v)=\delta(v)+Av\)，其中 \(\delta\) 逐坐标作用。另一方面，\(K\) 线性映射 \(f:M_L\to K\) 为 \(D\) 模同态，当且仅当其取值列 \(Y=(f(e_0),\ldots,f(e_{m-1}))^{\mathsf T}\) 满足
\[
\delta(Y)=A^{\mathsf T}Y.
\]
这些取值列恰为 \(Y=(y,\delta y,\ldots,\delta^{m-1}y)^{\mathsf T}\)，其中 \(L(y)=0\)。特别地，\(m=1\) 时 \(A=(-a_0)\)，系统为 \(\delta y=-a_0y\)。
\end{proposition}
\begin{proof}
由\cref{b2:prop:differential-cyclic-module}，所列 \(e_i\) 确为基。左乘 \(\partial\) 使 \(\partial e_j=e_{j+1}\)、\(j<m-1\)，而 \(Le_0=0\) 给出
\[
\partial e_{m-1}=-\sum_{i=0}^{m-1}a_ie_i.
\]
所以 \(A\) 的第 \(j\) 列正是 \(\partial e_j\) 的坐标。对 \(v=\sum_jv_je_j\)，Leibniz 公式给出
\[
\partial v=\sum_j\delta(v_j)e_j+\sum_jv_j\partial e_j,
\]
即坐标式 \(\delta(v)+Av\)。

若 \(f\) 为 \(D\) 模同态，则
\[
\delta\bigl(f(e_i)\bigr)=f(\partial e_i)
=\sum_j A_{ji}f(e_j),
\]
这就是 \(\delta(Y)=A^{\mathsf T}Y\)：模中基元素的作用记在 \(A\) 的各列，而同态对这些基元素的取值将同一列系数收集成系统的一行。反过来，若 \(Y\) 满足该系统，令 \(f(\sum_jv_je_j)=\sum_jv_jY_j\)，则
\[
f(\partial v)
=\sum_j\delta(v_j)Y_j+\sum_jv_j\delta(Y_j)
=\delta\bigl(f(v)\bigr).
\]
它已经保持 \(K\) 作用，又与 \(\partial\) 作用交换，而 \(D\) 由 \(K\) 和 \(\partial\) 生成，故为 \(D\) 模同态。

系统的前 \(m-1\) 行为 \(\delta(Y_i)=Y_{i+1}\)，逐次得到 \(Y_i=\delta^i(Y_0)\)；最后一行则是 \(\delta^m(Y_0)=-\sum_i a_i\delta^i(Y_0)\)，即 \(L(Y_0)=0\)。反向把任意这样的解及其连续导数代入，就满足全部系统。\(m=1\) 时没有前面的递推行，最后一行仍给出所述一阶方程。
\end{proof}

上面的转置来自模坐标与同态取值的不同约定，并不把 \(\nabla\) 变成 \(K\) 线性算子。回到一般 Ore 扩张，能否把系数放在右侧，同样取决于系数关系是否能够反向求解。

\begin{proposition}\label{b2:prop:ore-right-normal-form}
设 \(R\) 为含幺环，\(\sigma:R\to R\) 为保幺环自同构，\(\delta:R\to R\) 为加法映射，且对任意 \(a,b\in R\) 满足 \(\delta(ab)=\sigma(a)\delta(b)+\delta(a)b\)。则 Ore 扩张 \(R[t;\sigma,\delta]\) 作为右 \(R\) 模也以 \(1,t,t^2,\ldots\) 为基，每个元素唯一写成有限和 \(\sum_i t^ib_i\)、\(b_i\in R\)。

若只要求 \(\sigma\) 为保幺自同态，这个结论可能失败。具体地，设 \(k\) 为域，\(R=k\times k\)、\(\sigma(a,b)=(a,a)\)、\(\delta=0\)，令 \(e=(0,1)\)。则在 \(R[t;\sigma]\) 中 \(te=0\)、\(et\ne0\)，而 \(et\) 没有右系数表达 \(\sum_i t^ib_i\)。
\end{proposition}
\begin{proof}
将基本关系用于 \(\sigma^{-1}(a)\)，便得
\[
at=t\sigma^{-1}(a)-\delta\bigl(\sigma^{-1}(a)\bigr).
\]
这使一个系数能够从 \(t\) 的左侧移到右侧。进一步归纳可得
\[
at^n=t^n\sigma^{-n}(a)+\sum_{i<n}t^ib_i
\]
对适当的 \(b_i\in R\) 成立：将 \(at^{n+1}=(at^n)t\) 中每个右侧系数 \(b\) 按上述 \(bt\) 公式再移一次即可。因此每个左正规式都可改成右系数表达。

若有非平凡的有限关系 \(\sum_i t^ib_i=0\)，取最大的 \(n\) 使 \(b_n\ne0\)。由基本关系归纳，\(t^ib_i\) 的左系数表达最高项为 \(\sigma^i(b_i)t^i\)。故所给关系的 \(t^n\) 系数为 \(\sigma^n(b_n)\ne0\)，与\cref{b2:thm:ore-extension-construction}的左正规式唯一性矛盾。这证明右线性独立。

在第二种情形，\(\sigma\) 确为保幺环自同态，零映射满足扭曲导子公式。由 \(\sigma(e)=0\) 得 \(te=0\)，而左正规式定理给出 \(et\ne0\)。对 \(i\ge1\)，有
\[
t^i(a,b)=\sigma^i(a,b)t^i=(a,a)t^i.
\]
若 \(et=\sum_i t^ib_i\)，左正规式比较一次项就要求某个 \((a,a)\) 等于 \((0,1)\)，不可能。因此右侧表达的存在性失败；\(te=0\) 还直接说明这些幂不是右线性无关的。
\end{proof}

第三卷附录E继续介绍微分代数与消去；第五卷附录E讨论微分模及微分 Galois 理论，附录D则简介 D-模的几何观点。代数上的微分作用与Hilbert空间上的有界作用还有区别。\appref{b2:app:D}的\cref{b2:prop:bounded-weyl-impossible}证明，非零含幺 Banach 代数\footnote{巴拿赫（Stefan Banach，1892-1945），波兰数学家，建立赋范线性空间理论，著有《线性算子理论》。}中不能出现交换子等于一的两个元素。因此研究这些微分算子的分析表示时，必须另外考察无界算子及其定义域。

% !TeX root = ../../Book2.tex
% 纲要标题：### 20.6* 非交换 Gröbner 基与 Diamond 引理
% 纲要位置：计划纲要.md，第 1731 行。

% 纲要要点（供目录核对）：
%
% * 自由结合代数中的可容许字序、重写规则及包含和重叠歧义；下降次序须与两侧连接相容，两个最小失败例的实际商在正文命题完整证明
% * Diamond 引理与正规字；分别证明终止、首步选择无关及正规映射核等于关系理想，不把交换 Buchberger 判据原样照搬
% * PBW 型代数及可解型多项式代数中的有效约化；区别单次约化终止、全部歧义可解、规则有效匹配与补全终止，有限展示不自动给出有限 Gröbner 基
% * 建议先读第一卷附录 F.1—F.3 的交换 Gröbner 基；第三卷第8章发展计算应用、附录 F 处理局部标准基；正文用同一理想的两种字序证明有限基依赖字序，原题改为乘法、中心与单性计算，不把维数数据当作乘法结构
% ---
%

\OptionalSection{非交换 Gröbner 基与 Diamond 引理}{b2:sec:2006}

第一卷附录F.1—F.3讨论交换多项式的除法、首项理想和Buchberger 算法\footnote{布赫贝格尔（Bruno Buchberger，1942-），奥地利数学家，提出 Gröbner 基算法。}。本节在这些概念的背景下重新定义非交换字的次序与重写规则，说明字母顺序带来的差别。

\ProseParagraphBreak
正规式不仅要能够算出，还须证明不同计算次序给出相同结果。本节在自由结合代数的字上建立一个首一重写版本的 Diamond 引理。规则下降保证每次约化过程终止，歧义可解才保证不同选择得到相同正规式。要把它用于算法，还须说明规则如何匹配、系数如何计算；有限条原始关系也不保证补全后仍只有有限条规则。这些是不同的要求。


\subsection{可容许字序、重写规则与歧义}
\label{b2:subsec:2006:01}

令 \(X\) 为有限字母集，\(X^*\) 为全部有限字，空字记为一。自由结合代数 \(k\langle X\rangle\) 以这些字为基，乘法为连接。

\begin{definition}\label{b2:def:admissible-word-order}
设 \(X\) 为有限字母集，\(X^*\) 为由 \(X\) 构成的全部有限字的集合，空字记为 \(1\)。若 \(X^*\) 上的全序 \(<\) 为良序，以空字为最小元，且对任意字 \(a,b,u,v\in X^*\)，\(u<v\) 蕴含 \(aub<avb\)，则称 \(<\) 为\term[kerongxu-zixu]{可容许字序}[admissible word order]。例如先比较字长，再按固定的字母次序作字典比较，就满足这些条件。
\end{definition}

对一个非零多项式 \(f\)，其最大的字称为\term[shou-danxiangshi]{首单项式}[leading monomial]，记为 \(\operatorname{LM}(f)\)。若其系数为一，称 \(f\) 首一。规定一族规则 \(\ell_\alpha\to r_\alpha\)，其中 \(\ell_\alpha\ne1\)，右侧为有限线性组合且其中每个字都严格小于 \(\ell_\alpha\)。这相当于使用首一关系 \(g_\alpha=\ell_\alpha-r_\alpha\)。

\ProseParagraphBreak
一次重写将多项式中一个完整项 \(c\,a\ell_\alpha b\) 换成 \(c\,ar_\alpha b\)，然后合并同类项。字若不含任何 \(\ell_\alpha\) 作为连续子字，称为\term[zhenggui-zi]{正规字}[normal word]。右侧的字小于 \(\ell_\alpha\)，还必须在加入任意前缀 \(a\)、后缀 \(b\) 后继续小于 \(a\ell_\alpha b\)。可容许字序的双侧相容性正保证这件事；仅有良序而没有这项相容性，局部的下降不能保证大字内部的重写也下降。

\ProseParagraphBreak
同一字中若有两处可替换，可能出现两类\term[chongxie-qiyi]{歧义}[rewriting ambiguity]。一类是包含：\(\ell_\alpha=a\ell_\beta b\)，可先改整个左端，也可先改其中一段；另一类是重叠：\(\ell_\alpha=au\)、\(\ell_\beta=ub\)，其中 \(a,u,b\) 都非空，字 \(aub\) 有两种第一步。相同左端的不同规则也算包含歧义。若两种第一步结果能继续重写为同一个多项式，称这项歧义可解。

\ProseParagraphBreak
例如规则 \(x^2\to1\)、\(yx\to0\) 在 \(yxx\) 上重叠：改末尾 \(xx\) 得到 \(y\)，改开头 \(yx\) 却得到零。规则 \(xyx\to x\)、\(yx\to0\) 则在 \(xyx\) 上包含：改整个字得到 \(x\)，改其中的末尾 \(yx\) 得到零。两组原始规则都下降，但各自的两支已不能继续合到一起。\cref{b2:prop:diamond-ambiguity-failures}将计算这两组关系实际定义的商代数。

\ProseParagraphBreak
两处彼此分离时不形成新的歧义：在字 \(a\ell_\alpha b\ell_\beta c\) 中，先改哪一处，再改另一处，都会得到 \(ar_\alpha br_\beta c\)。这里利用有限和的分配律，不交换任何字母。

\subsection{Diamond 引理与正规字}
\label{b2:subsec:2006:02}

\begin{theorem}[Diamond 引理，首一字重写版本]\label{b2:thm:diamond-lemma}
设 \(k\) 为域，\(X\) 为有限字母集，\(<\) 为 \(X^*\) 上的可容许字序。给定一族首一字重写规则 \(\ell_\alpha\rightarrow r_\alpha\)，其中每个 \(\ell_\alpha\) 都是非空字，\(r_\alpha\) 是严格小于 \(\ell_\alpha\) 的字的有限 \(k\) 线性组合。每次重写在多项式中将一个完整项 \(c\,a\ell_\alpha b\) 换为 \(c\,ar_\alpha b\)，其中 \(c\in k\)，\(a,b\in X^*\)，然后合并同类项，则从任意多项式出发的每个重写过程都在有限步后停止。若全部包含与重叠歧义可解，则每个多项式有唯一的正规字线性组合形式；这里正规字指不含任何 \(\ell_\alpha\) 为连续子字的字。令 \(I\) 为全部 \(\ell_\alpha-r_\alpha\) 生成的双边理想，则正规字的商类构成 \(k\langle X\rangle/I\) 的一组 \(k\) 基。
\end{theorem}

这项结果称为\term[Diamond-yinli]{Diamond 引理}[Diamond lemma]。这里证明的是自由结合代数上采用首一规则与可容许全序的专门形式；原论文还讨论更一般的约化系统\cite{Bergman1978DiamondLemma}。只写若干等式而不给下降次序，还不能使用它。

\begin{proof}
每次重写将一个完整项换为有限多个更小的字，下降依据是良序，右侧的有限性则保证仍得到有限多项式。当前最大字未必每步都下降，因为也可以先改较小项；要证明任意选择都终止，对有限多项式的最大字 \(w\) 作良序归纳。只要 \(w\) 尚未被替换，对较小项进行的操作都是从一个有限的较小字多项式出发，按归纳假设不可能无限进行。较小字也不能产生 \(w\)。若 \(w\) 可替换，把它的整个系数项换掉后，余下所有字都小于 \(w\)，再次由归纳假设终止；若 \(w\) 不可替换，只需处理下面各项。最小字一没有可用规则，提供归纳起点。合并同类项造成的消去只会删除项，不会破坏这个下降论证。

再对单个字 \(w\) 作良序归纳证明正规结果唯一，并记为 \(N(w)\)。假设所有更小的字已经有唯一正规结果，将 \(N\) 线性延伸到它们的有限和上。若 \(w\) 正规，规定 \(N(w)=w\)。否则，选择一次替换得到多项式 \(f\)，其各字小于 \(w\)，候选结果为 \(N(f)\)。须证明它不依赖第一步。

比较两种第一步。如果两处互不相交，继续改另一处会得到上一小节写出的同一个表达式，因此按归纳假设两者的正规结果相同。如果两处包含或重叠，它们在某个连续子字中构成一项题设歧义。该歧义的两支能化为共同结果；在两侧加上原字的相同前缀和后缀，仍是合法的重写序列。所有这些字在第一步以后都严格小于 \(w\)，所以归纳假设再次保证两支的正规结果相同。这就排除了第一步的全部选择，并完成归纳。

由线性延伸，得到到正规字线性生成空间的映射 \(N\)。每次重写都保持 \(N\)，所以每个终止结果等于 \(N(f)\)。又每次重写之差属于 \(I\)，因而 \(f-N(f)\in I\)，这已经保证正规字的商类生成商代数。为证明独立，还要排除正规字之间藏在 \(I\) 中的非零关系。对任意字 \(a,b\) 和规则，有 \(N(a\ell_\alpha b)=N(ar_\alpha b)\)，所以 \(N\) 零化 \(a(\ell_\alpha-r_\alpha)b\)。这些元素线性生成 \(I\)，故 \(N(I)=0\)，即 \(I\subseteq\ker N\)；而 \(N(f)=0\) 时，\(f=f-N(f)\in I\)，给出反向包含。因此 \(\ker N=I\)。正规字上的 \(N\) 是恒等映射，所以它们的非零线性组合不可能属于 \(I\)，得到独立性。
\end{proof}

例如选择 \(y>x\) 的长度字典序，规则 \(yx\to xy+1\) 没有自身的包含或重叠歧义：左端的非空真后缀为 \(x\)，非空真前缀为 \(y\)，它们不同。正规字恰为 \(x^iy^j\)，所以引理给出 Weyl 代数的正规单项式基的另一证明。规则 \(yx\to qxy\) 同样处理量子平面；这里依赖的是字序和歧义检查。

\subsection{PBW 型代数的有效约化}
\label{b2:subsec:2006:03}

当规则有限且系数可精确计算时，可以固定一种约化程序：每步取最大的可约字，再选取其中最左的可约位置；若仍有多条规则，按事先编号选择。上一定理保证停止。若歧义条件也成立，输出就与这个具体选择无关，并且
\[
f\in I\quad\Longleftrightarrow\quad N(f)=0.
\]
因此正规式可以用于相等判定与理想成员判定。终止性本身没有给出统一的时间复杂度界。

\ProseParagraphBreak
例如 \(yx\to xy+x^2\) 在 \(y>x\) 的长度字典序下严格下降，且仍没有自身歧义，因此 \(x^iy^j\) 也构成正规单项式基。计算得
\[
yx^2=x^2y+2x^3,\qquad
y^2x=xy^2+2x^2y+2x^3.
\]
这个代数也可用\cref{b2:thm:ore-extension-construction}构造：取 \(R=k[x]\)、\(\sigma=\operatorname{id}\)、\(\delta(f)=x^2f'\)，再把扩张变量记为 \(y\)。两种方法一个检查重写歧义，一个通过导子建立乘法，得到同一条关系及同一组基。

\begin{definition}\label{b2:def:noncommutative-groebner-basis}
设 \(k\) 为域，\(X\) 为有限字母集，\(<\) 为 \(X^*\) 上的可容许字序，\(I\subseteq k\langle X\rangle\) 为双边理想，\(G\subseteq I\setminus\{0\}\) 为 \(I\) 的一组生成元。若每个非零 \(f\in I\) 的首单项式都含某个 \(g\in G\) 的首单项式为连续子字，则称 \(G\) 为 \(I\) 关于 \(<\) 的\term[feijiaohuan-Groebner-ji]{非交换 Gröbner 基}[noncommutative Gröbner basis]。
\end{definition}

满足 Diamond 引理条件的首一关系族就是这样的基。事实上，若 \(f\in I\) 的最大字不可约，它在任何对较小字的重写中都不可能被消去，所以 \(N(f)\ne0\)，与 \(N(I)=0\) 矛盾。

\ProseParagraphBreak
若某项歧义产生两个不同的正规结果，它们之差是理想中的新关系；可以将差的首项归一化，加入新的下降规则，再检查新出现的歧义。这种补全同样利用尚未消去的关系来扩大首项条件，但非交换字的歧义来自连续包含与连续重叠，不能用交换单项式的最小公倍数替代。即使原规则有限，每次新增的左端仍可能产生下一项歧义，不能由单次约化会停止推断整个补全过程也会停止。

\subsection{非交换 Gröbner 基的有限性问题}
\label{b2:subsec:2006:04}

有限条原始关系不保证只需补充有限多条规则，甚至不保证在给定字序下存在任何有限 Gröbner 基。下面对一个仅有两个生成元、一条关系的代数证明这一点：不仅所写的补全过程无穷延续，其首字反链还将排除所有有限替代。不过这个结论依赖指定的字序，换序后同一理想可以有有限 Gröbner 基。

\ProseParagraphBreak
在 \(k\langle x,y\rangle\) 中，取 \(x>y\) 的长度字典序，从原关系 \(x^2-xy\) 得到规则 \(x^2\to xy\)。对同一个字 \(x^3\)，两种第一步产生分叉：
\[
(x^2)x\longrightarrow xyx,
\qquad
x(x^2)\longrightarrow x^2y\longrightarrow xy^2.
\]
两支的末项都已不能用原规则继续约化，却应代表同一个商类。为使这项歧义可解，须补上 \(xyx\to xy^2\)。

\ProseParagraphBreak
这种补充会继续下去。若已补到规则 \(xy^nx\to xy^{n+1}\)，它与原规则在 \(xy^nx^2\) 上又有两种约化：
\[
\begin{aligned}
(xy^nx)x&\longrightarrow xy^{n+1}x,\\
xy^n(x^2)&\longrightarrow xy^nxy\longrightarrow xy^{n+2}.
\end{aligned}
\]
于是下一条规则应是 \(xy^{n+1}x\to xy^{n+2}\)。这些左端中的两个 \(x\) 相隔越来越多的 \(y\)，以前的规则不能直接改写新的左端。下面证明这一过程产生的全部关系确实属于原理想，并形成无限 Gröbner 基。

\begin{proposition}\label{b2:exm:infinite-noncommutative-groebner}
设 \(k\) 为域，在 \(k\langle x,y\rangle\) 中令 \(I=(x^2-xy)\)。采用 \(x>y\) 的长度字典序时，该理想没有有限非交换 Gröbner 基。一组无限首一 Gröbner 基为
\[
g_n=xy^nx-xy^{n+1},\qquad n\ge0.
\]
商代数的正规字为 \(y^a\) 及 \(y^axy^b\)、\(a,b\ge0\)。改用 \(y>x\) 的长度字典序时，单条关系 \(xy-x^2\) 就是有限首一 Gröbner 基，相应正规字为 \(y^ax^b\)、\(a,b\ge0\)。
\end{proposition}
\begin{proof}
\(g_0=x^2-xy\) 是原关系，而直接展开给出
\[
g_{n+1}=g_ny-g_nx+xy^n g_0.
\]
由归纳，全部 \(g_n\in I\)，所以它们生成的双边理想仍恰为 \(I\)。规则 \(xy^nx\to xy^{n+1}\) 保持字长，最后一个 \(x\) 改成较小的 \(y\)，因此下降。

不同左端不互相包含，因为它们各有且只有两个 \(x\)。非平凡重叠只能共享这个末端与另一个左端的首个 \(x\)，得到字 \(xy^axy^bx\)。先改左边，得到 \(xy^{a+b+1}x\)，再改为 \(xy^{a+b+2}\)；先改右边，得到 \(xy^axy^{b+1}\)，再改左边，也得到 \(xy^{a+b+2}\)。因此所有歧义可解。\cref{b2:thm:diamond-lemma}给出正规字基，而不可约字恰好是至多含一个 \(x\) 的字，即所列形式。

这也确定了理想的全部首单项式：它们恰为包含某个 \(xy^nx\) 的字。一个方向由每个 \(g_n\in I\) 及两边乘字得到；另一个方向由正规式判定得到，因为理想中的非零元素不可能有不可约的最大字。

各个字 \(xy^nx\) 对连续子字包含关系构成无限反链，因而理想的这些最小首字不能由有限多个其他首字统一覆盖。具体地，若存在有限 Gröbner 基，对每个 \(n\)，其中某个首单项式必须是 \(xy^nx\) 的子字；而它又必须包含某个 \(xy^mx\)。因为两个 \(x\) 必须同时保留，只能有 \(m=n\)，且该首单项式就是整个 \(xy^nx\)。于是基中必须为每个 \(n\) 提供一个不同的首单项式，与有限性矛盾。

改用 \(y>x\) 的长度字典序后，\(xy-x^2\) 的首字为 \(xy\)，规则 \(xy\to x^2\) 严格下降。它没有非平凡的自身包含歧义，也没有重叠歧义，因为唯一的非空真后缀 \(y\) 与唯一的非空真前缀 \(x\) 不同。因此\cref{b2:thm:diamond-lemma}给出正规字基，并由前面首项判定说明这条关系构成 Gröbner 基。不含连续子字 \(xy\) 的字只能先写若干个 \(y\)，再写若干个 \(x\)，恰为 \(y^ax^b\)。这证明同一理想的有限 Gröbner 基存在性确实可以随字序改变。
\end{proof}

上述理想在 \(x>y\) 字序下的无限规则仍可用于计算：给定长度为 \(N\) 的字，只可能匹配 \(0\le n\le N-2\) 的左端 \(xy^nx\)。沿字从左到右寻找相邻两个 \(x\) 之间全由 \(y\) 组成的连续子字，就能判定可约位置并执行下一步；对有限多项式逐项检查即可。每步都在可容许字序下下降，所以在系数运算有效时得到终止的正规式计算。“没有有限 Gröbner 基”不等于“任何计算都不可能”。这里证明有限展示不保证给定字序下有有限 Gröbner 基，同时另行证明了这组无限规则可以有效匹配；不能仅由有限展示推得后一性质。第一卷附录F已经给出交换 Gröbner 基的有限性及基本算法；第三卷第8章进一步发展交换情形的计算与几何应用，第三卷附录F则处理局部标准基。使用时须分别核对对象、字序或项序，以及终止性的实际依据。

\begin{proposition}\label{b2:prop:diamond-ambiguity-failures}
设 \(k\) 为域。在 \(k\langle x,y\rangle\) 中，规则 \(x^2\to1\)、\(yx\to0\) 都按长度字典序下降，却在 \(yxx\) 上有不可解的重叠歧义，且
\[
k\langle x,y\rangle/(x^2-1,yx)\cong k[s]/(s^2-1).
\]
在这个商中 \(y=0\)，基为 \(1,x\) 的商类。规则 \(xyx\to x\)、\(yx\to0\) 也都下降，却在 \(xyx\) 上有不可解的包含歧义，且
\[
k\langle x,y\rangle/(xyx-x,yx)\cong k[s].
\]
在这个商中 \(x=0\)，基为 \(1,y,y^2,\ldots\) 的商类。
\end{proposition}
\begin{proof}
第一组规则在 \(yxx\) 上，一条约化先改 \(xx\) 而得到 \(y\)，另一条先改 \(yx\) 而得到零；\(y\) 与零对于原规则都已正规，所以歧义不可解。令 \(A=k\langle x,y\rangle/(x^2-1,yx)\)，暂将生成元的商类仍记为 \(x,y\)。结合律与两条关系给出
\[
y=yx^2=(yx)x=0.
\]
向 \(k[s]/(s^2-1)\) 发送 \(x\mapsto s\)、\(y\mapsto0\)，两条关系均变成零，故得到同态 \(A\to k[s]/(s^2-1)\)。反向以 \(s\mapsto x\) 定义同态，因 \(x^2=1\) 而良定义。两个复合保持生成元；在 \(A\) 中，已经证明 \(y=0\)，所以这两个同态互逆。首一多项式 \(s^2-1\) 的带余除法给出 \(1,s\) 的商类为基，得到 \(A\) 中所列的基。这个论证在特征二时也成立。

第二组规则中，改整个 \(xyx\) 得到 \(x\)，改末尾 \(yx\) 却得到零；\(x\) 与零对原规则仍都正规。令 \(B=k\langle x,y\rangle/(xyx-x,yx)\)，其关系给出
\[
x=xyx=x(yx)=0.
\]
发送 \(x\mapsto0\)、\(y\mapsto s\) 得到同态 \(B\to k[s]\)，反向以 \(s\mapsto y\) 得到同态 \(k[s]\to B\)。两个复合都保持生成元，其中 \(B\) 内的 \(x=0\)，所以它们互逆。\(k[s]\) 的幂基给出所列商类为基。两组失败都产生了原先终止结果未显示的额外关系，因此原规则的终止性没有保证正规字的商类线性独立。
\end{proof}


\begin{chaptersummary}
滤过保留“不超过给定次数”的信息，伴随分次取每一层的新部分。穷尽性保证元素都有有限滤过次数，下界使逐次消去主符号的过程停止；由伴随分次提升 Noether 性时，子模须用诱导滤过。Rees 代数的零纤维为伴随分次，非零纤维恢复原代数，但平坦性不要求纤维具有相同乘法或相同零因子性质。粗化的多项式滤过产生平方零主符号，已经表明原代数为整环不保证伴随分次也是整环。

\ProseParagraphBreak
有序单项式的生成性与独立性分别需要证明。Weyl 关系的重排给出生成性，辅助变量作用 \(x^i\mathrm d^j(1)=X^iY^j\) 给出独立性；量子平面的扭曲乘法直接构造具有独立基的代数；Ore 构造用 \(T^i(1)=z^i\) 读取左正规式的系数。右正规式还须另查 \(\sigma\) 的可逆性。Weyl 代数的中心在特征 \(p\) 时扩大，中心上的秩为 \(p^2\)，而量子平面的中心由参数的精确阶决定，有限阶 \(\ell\) 时相应秩为 \(\ell^2\)。特征零的 Weyl 代数无非零有限维幺性表示；正特征则可由固定中心值的商得到矩阵代数。

\ProseParagraphBreak
Lie 括号的 Jacobi 恒等式保证伴随作用是表示，泛包络代数把括号关系转为结合乘法，使 Lie 表示对应于幺性左模。Heisenberg 的中心作用为零时，参数商是交换多项式环；作用为非零标量时，商经缩放成为 Weyl 代数。一般 PBW 定理确定 \(\operatorname{gr}U(\mathfrak g)\cong\operatorname{Sym}(\mathfrak g)\)，从而保证有序单项式的独立性。仅把写出的关系取最高次部分，不能代替这个同构的证明：坏 PBW 关系的 \(zyx\) 分叉与 Jacobi 失效产生同一个隐藏关系 \(x=0\)，实际伴随分次也会记录这一坍缩；剩余代数仍可构造正规基。

\ProseParagraphBreak
线性微分方程由左理想商 \(M_L=D/DL\) 编码，函数解是 \(\operatorname{Hom}_D(M_L,K)\) 的元素。模元素的坐标作用要保留系数求导项，模同态的取值系统则使用转置矩阵。Diamond 引理中，下降字序保证约化终止，歧义可解保证正规结果唯一；补充规则是否终止是另一问题。理想 \((x^2-xy)\) 在 \(x>y\) 的字序下无有限 Gröbner 基，换成 \(y>x\) 时一条规则已足够。商的每次分次维数虽与二元多项式环相同，仍有 \(x(x-y)=0\)：维数数据无法单独恢复乘法结构。
\end{chaptersummary}

\BookExercises

\begin{optionalexercise}\label{b2:ex:20-coarse-filtration}
在 \(A=k[x]\) 上取 \(F_nA=\{f\mid\deg f\le2n\}\)，负次滤过项为零。证明其 Rees 代数有展示
\[
\mathcal R_F(A)\cong k[Z,U,V]/(U^2-ZV),
\qquad Z\mapsto z,\quad U\mapsto xz,\quad V\mapsto x^2z.
\]
直接计算 \(Z=0\) 与 \(Z=1\) 的两个商，并解释平方零元为何只出现在第一个商中。
\end{optionalexercise}
\begin{solution}
每个 Rees 单项式 \(x^iz^n\) 满足 \(i\le2n\)。写 \(i=2b+a\)、\(a\in\{0,1\}\)，则 \(n\ge b+a\)，所以它等于 \(U^aV^bZ^{n-b-a}\) 的像。关系 \(U^2=ZV\) 把每个交换单项式化为 \(Z^cU^aV^b\)、\(a=0,1\)。这些元素的像是 \(x^{a+2b}z^{a+b+c}\)；不同三元组给出不同单项式，因为 \(x\) 次数先确定 \(a,b\)，再由 \(z\) 次数确定 \(c\)。故所给满同态也单射。

模 \(Z\) 后得到 \(k[U,V]/(U^2)\)，其中 \(U\ne0\) 而 \(U^2=0\)；它就是\cref{b2:prop:filtered-domain-converse-failure}中的伴随分次。模 \(Z-1\) 后关系成为 \(V=U^2\)，得到 \(k[U]\cong k[x]\)。原 Rees 代数中 \(U^2=ZV\ne0\)，取零纤维才把右侧杀掉；因此平坦族的不同纤维可以有不同的零因子性质。
\end{solution}

\begin{optionalexercise}\label{b2:ex:20-polynomial-rees}
在 \(k[x,y]\) 中赋予 \(x,y\) 权重一、二，以权重不超过 \(n\) 的单项式张成 \(F_nA\)。证明
\[
\mathcal R_F(A)=k[z,xz,yz^2]\cong k[Z,U,V],
\qquad \deg Z=\deg U=1,\quad\deg V=2.
\]
辨认零纤维与一纤维，并说明零纤维的分次如何记录原来两个变量的不同权重。
\end{optionalexercise}
\begin{solution}
Rees 单项式为 \(x^iy^jz^n\)、\(i+2j\le n\)，故等于 \((xz)^i(yz^2)^jz^{n-i-2j}\)。不同的 \(Z^aU^iV^j\) 映到不同的 \(x^iy^jz^{a+i+2j}\)，这给出所述多项式环同构和次数。模 \(Z\) 后得到分次环 \(k[U,V]\)，其中 \(U,V\) 分别对应 \(x\) 的一次主符号与 \(y\) 的二次主符号；模 \(Z-1\) 后得到原环 \(k[x,y]\)。两纤维作为未分次的环都是二元多项式环，但零纤维保留的权重分次并不是两个变量都为一次的通常分次。
\end{solution}

\begin{optionalexercise}\label{b2:ex:20-filtered-isomorphism}
设 \(k\) 为域，\(M,N\) 为具有非负、穷尽滤过的 \(k\) 向量空间，\(k\) 线性映射 \(f:M\to N\) 满足 \(f(F_nM)\subseteq F_nN\)。若 \(\operatorname{gr}f\) 是同构，证明 \(f\) 是同构，且 \(f(F_nM)=F_nN\) 对每个 \(n\) 成立。
\end{optionalexercise}
\begin{solution}
若 \(0\ne m\in\ker f\)，取其最低滤过次数，则它的非零主符号被 \(\operatorname{gr}f\) 送到零，矛盾。故 \(f\) 单射。对 \(y\in F_nN\)，商片上的满射性给出 \(x\in F_nM\)，使 \(y-f(x)\in F_{n-1}N\)。归纳提升这个差值，在 \(F_{-1}N=0\) 处结束，就得 \(y\in f(F_nM)\)。穷尽性再使 \(f\) 满射。整个论证同时证明每一滤过项上的满射性。
\end{solution}

\begin{optionalexercise}\label{b2:ex:20-weyl-normal-ordering}
在任意特征的 Weyl 代数中证明
\[
\mathrm{d}^mx^n=\sum_{r=0}^{\min(m,n)}\binom mr n(n-1)\cdots(n-r+1)x^{n-r}\mathrm{d}^{m-r},
\]
其中 \(r=0\) 的下降乘积为一，整数系数取其在 \(k\) 中的像。
\end{optionalexercise}
\begin{solution}
\(m=0\) 时成立。由 \(\mathrm{d}f(x)=f(x)\mathrm{d}+f'(x)\)，对第 \(m\) 步的每项再左乘 \(\mathrm{d}\)，得到保持导数次数的一项及增加一次导数的一项。相同 \(r\) 的系数由 \(\binom mr+\binom m{r-1}=\binom{m+1}r\) 合并，下降乘积的新增因子为 \(n-r+1\)。这正给出第 \(m+1\) 步。证明只用整数恒等式及形式导数，没有把可能为零的阶乘作除数，故对任意特征成立。
\end{solution}

\begin{optionalexercise}\label{b2:ex:20-weyl-finite-representations}
设 \(\operatorname{char}k=p>0\)，\(V\ne0\) 为有限维幺性左 \(A_1(k)\) 模，且 \(x^pV=\mathrm d^pV=0\)。证明 \(V\) 是 \(k[X]/(X^p)\) 的若干个直和；说明直和项数如何由 \(\dim_kV\) 决定，以及 \(V\) 何时简单。
\end{optionalexercise}
\begin{solution}
由\cref{b2:prop:weyl-finite-dimensional-representations}，作用经过 \(A_1(k)/(x^p,\mathrm d^p)\cong M_p(k)\)。记这个矩阵环的矩阵单位为 \(E_{ij}\)，其中 \(0\le i,j<p\)，令 \(W=E_{00}V\)。恒等式 \(1=\sum_iE_{ii}\) 给出 \(V=\bigoplus_iE_{ii}V\)，而 \(E_{i0}:W\to E_{ii}V\) 与 \(E_{0i}\) 互逆。因此
\[
k^p\otimes_kW\longrightarrow V,\qquad e_i\otimes w\longmapsto E_{i0}w
\]
是 \(M_p(k)\) 模同构，矩阵环只作用于第一个因子。正文的截断多项式表示正是标准列模 \(k^p\)。取 \(W\) 的一组基，得到 \(V\cong\bigl(k[X]/(X^p)\bigr)^{\oplus r}\)，其中 \(r=\dim_kW=\dim_kV/p\)。标准列模简单；若 \(r>1\)，一个直和项给出真非零子模，所以 \(V\) 简单当且仅当 \(r=1\)。
\end{solution}

\begin{optionalexercise}\label{b2:ex:20-weyl-characteristic-p-matrix}
设 \(\operatorname{char}k=p>0\)，\(a,b\in k\)。证明
\[
A_1(k)/(x^p-a^p,\mathrm d^p-b^p)\cong M_p(k).
\]
在 \(k[X]/(X^p)\) 上写出这个商的简单表示，并求其作为 \(A_1(k)\) 表示的核。
\end{optionalexercise}
\begin{solution}
置 \(u=x-a\)、\(v=\mathrm d-b\)，则 \(vu-uv=1\)。标量与生成元交换，故特征 \(p\) 的二项式公式给出 \(u^p=x^p-a^p\)、\(v^p=\mathrm d^p-b^p\)。代换与逆代换 \(x=u+a,\mathrm d=v+b\) 因而把所给商同构到 \(A_1(k)/(u^p,v^p)\)，由\cref{b2:prop:weyl-finite-dimensional-representations}它是 \(M_p(k)\)。

相应作用为 \(x\cdot f=(X+a)f\)、\(\mathrm d\cdot f=f'+bf\)。减去标量后恢复正文的截断多项式表示，因此这个表示简单。商矩阵环在标准列模上的作用忠实，故原代数的作用核恰为双边理想 \((x^p-a^p,\mathrm d^p-b^p)\)。这里只处理中心值已经是 \(k\) 中 \(p\) 次幂的纤维，未假定任意中心值都能作这样的变量平移。
\end{solution}

\begin{optionalexercise}\label{b2:ex:20-weyl-center-module-rank}
设 \(\operatorname{char}k=p>0\)，\(C=Z(A_1(k))\)、\(F=\operatorname{Frac}(C)\)。证明 \(F\otimes_CA_1(k)\) 是中心为 \(F\)、维数为 \(p^2\) 的除代数，因而不是 \(M_p(F)\)。
\end{optionalexercise}
\begin{solution}
由\cref{b2:thm:weyl-center-simplicity}，\(C=k[x^p,\mathrm d^p]\)，且 \(A_1(k)\) 在 \(C\) 上以 \(x^i\mathrm d^j\)、\(0\le i,j<p\)，为自由基。原代数因而嵌入中心局部化；该局部化以同一组元素为 \(F\) 基，维数为 \(p^2\)。原 Weyl 代数是整环；中心非零元作分母后仍是整环，因为两个非零分式的乘积有非零分子。

有限维域代数中，非零元素的左乘在整环条件下单射，故满射，给出右逆；整环条件又使该右逆同时为左逆。因此这个局部化是除代数。若 \(c=\sum_{i,j<p}c_{ij}x^i\mathrm d^j\) 居中，所有系数属于 \(F\)，与 \(\mathrm d\) 取交换子使所有 \(i>0\) 的系数为零，与 \(x\) 取交换子再使所有 \(j>0\) 的系数为零；这里 \(1,\ldots,p-1\) 在 \(k\) 中都非零。故中心恰为 \(F\)。而 \(p\ge2\)，\(M_p(F)\) 有非零零因子，所以不能与这个除代数同构。有限秩和相同维数本身不决定中心上的代数是否分裂。
\end{solution}

\begin{optionalexercise}\label{b2:ex:20-quantum-units-ideals}
求量子平面 \(k_q[x,y]\) 的单位，并证明它对任意 \(q\ne0\) 都不是单环。特别解释中心等于 \(k\) 为什么不足以推出单性。
\end{optionalexercise}
\begin{solution}
由\cref{b2:prop:quantum-plane-domain}，非零乘积的 \(y\) 次数相加。若两个元素互逆，它们都不含 \(y\)，于是属于 \(k[x]\)，只能是非零常数。因此单位群为 \(k^\times\)。关系使 \(Ax=xA\)，双边理想 \((x)\) 由 \(x\) 指数为正的正规单项式生成，非零且不含一；商为 \(k[y]\)。所以环不单。参数无限阶时中心虽为 \(k\)，这仍不改变该真双边理想的存在。
\end{solution}

\begin{optionalexercise}\label{b2:ex:20-quantum-binomial}
在量子平面中把 \((x+y)^3\) 写成正规式。若 \(q\ne1\) 且 \(1+q+q^2=0\)，证明 \((x+y)^3\) 属于中心。
\end{optionalexercise}
\begin{solution}
先有 \((x+y)^2=x^2+(1+q)xy+y^2\)。再乘 \(x+y\)，用 \(yx=qxy\)、\(y^2x=q^2xy^2\)，得到
\[
(x+y)^3=x^3+(1+q+q^2)x^2y+(1+q+q^2)xy^2+y^3.
\]
题设使 \(q^3=1\) 且 \(q\ne1\)，所以参数的阶为三，式子化为 \(x^3+y^3\)。由\cref{b2:prop:quantum-plane-center}，两项均居中。
\end{solution}

\begin{optionalexercise}\label{b2:ex:20-quantum-inverse-parameter}
证明交换两个生成元给出 \(k_q[x,y]\cong k_{q^{-1}}[x,y]\)。说明它与中心阶数公式相容，但不要据此断言同阶的任意两个参数都给出同构代数。
\end{optionalexercise}
\begin{solution}
在目标代数中记生成元为 \(X,Y\)，有 \(YX=q^{-1}XY\)。赋值 \(x\mapsto Y,y\mapsto X\) 将原关系送到 \(XY-qYX=0\)，所以诱导同态。反向仍交换两生成元，两复合固定生成元，因而互逆。\(q\) 与 \(q^{-1}\) 的阶相同，故中心中的幂指数一致；这个构造只比较互逆的参数，没有处理其他同阶参数。
\end{solution}

\begin{optionalexercise}\label{b2:ex:20-weyl-rees-fibres}
固定正整数 \(r,s\)，在 Weyl 代数中令 \(x,\mathrm d\) 的权重分别为 \(r,s\)，以正规单项式 \(x^i\mathrm d^j\) 的权重 \(ri+sj\) 定义滤过。证明其 Rees 代数由 \(Z=z,X=xz^r,D=\mathrm dz^s\) 生成，且有展示
\[
k\langle X,D,Z\rangle/(ZX-XZ,ZD-DZ,DX-XD-Z^{r+s}).
\]
证明 \(X^iD^j\) 构成它在 \(k[Z]\) 上的一组基，并辨认各个 \(Z=\lambda\) 的纤维。
\end{optionalexercise}
\begin{solution}
Weyl 重排中的每个交换子项都将权重降低 \(r+s\)，故所给空间确为乘法相容的滤过。Rees 单项式 \(x^i\mathrm d^jz^n\)、\(ri+sj\le n\)，是 \(X^iD^jZ^{n-ri-sj}\) 的像。关系的像为零，反复移动 \(Z\) 并使用 \(DX=XD+Z^{r+s}\)，使任意字成为 \(X^iD^jZ^a\) 的线性组合：按 \(X,D\) 的字长及逆序数作下降归纳即可。它们的像 \(x^i\mathrm d^jz^{ri+sj+a}\) 由 Weyl 正规基及 \(z\) 次数线性无关，故没有其他关系，并得 \(k[Z]\) 自由基。

取纤维后仍以 \(X^iD^j\) 为 \(k\) 基，关系变为 \(DX-XD=\lambda^{r+s}\)。\(\lambda=0\) 时是交换多项式环；\(\lambda\ne0\) 时把 Weyl 的 \(x,\mathrm d\) 分别送到 \(X,\lambda^{-(r+s)}D\)，逆向将 \(X,D\) 送到 \(x,\lambda^{r+s}\mathrm d\)，得到互逆同构。指数 \(r+s\) 来自原关系最高权重，而不是参数的一次幂。
\end{solution}

\begin{optionalexercise}\label{b2:ex:20-bad-pbw-quotient}
固定整数 \(r\ge2\)，令 \(B_r=k\langle y,z\rangle/(zy-yz-y^r)\)。证明 \(B_r\cong k[y][z;\delta]\)，其中 \(\delta(f)=y^rf'\)，并给出正规单项式基。再令 \(y,z\) 的权重分别为一、\(r\)，证明相应滤过的伴随分次是加权多项式环。
\end{optionalexercise}
\begin{solution}
\(\delta\) 是 \(k[y]\) 的导子；Ore 构造给出 \(zy=yz+y^r\) 和基 \(y^iz^j\)。反向，从这个生成元关系归纳得到 \(zf=fz+y^rf'\)，再用 Ore 泛性质得到回到 \(B_r\) 的同态。两复合固定生成元，所以互逆，基结论随之成立。

赋权后，关系中的 \(zy,yz\) 权重为 \(r+1\)，\(y^r\) 权重为 \(r\)。移动每个逆序对时，交换项保持权重，导子项降低权重。因此权重不超过 \(n\) 的正规单项式恰好张成商上的第 \(n\) 个滤过项，且乘法相容。正规基的独立性使每个商片仍由权重恰为 \(n\) 的单项式组成；最高项关系只要求两个主符号交换，得到 \(\operatorname{gr}B_r\cong k[Y,Z]\)，其中 \(\deg Y=1,\deg Z=r\)。选择权重可将高次修正项变为低次项，但正规基的独立性仍须先由 Ore 构造保证。
\end{solution}

\begin{optionalexercise}\label{b2:ex:20-affine-lie-ideals}
设二维 Lie 代数 \(\mathfrak a=kh\oplus ke\) 满足 \([h,e]=e\)。求其全部 Lie 理想，并求全部一维表示。说明任何一维表示为什么都不能忠实。
\end{optionalexercise}
\begin{solution}
\(0,ke,\mathfrak a\) 显然都是理想。若理想 \(I\) 含有 \(ah+be\) 且 \(a\ne0\)，则 \([ah+be,e]=ae\in I\)，故 \(e\in I\)，继而 \(h\in I\)，于是 \(I=\mathfrak a\)。否则 \(I\subseteq ke\)，只有零空间与 \(ke\) 两种可能。

在一维空间上，每个算子为标量，算子交换子恒为零。表示条件因此给出 \(\rho(e)=\rho\bigl([h,e]\bigr)=0\)，而 \(\rho(h)=\lambda\) 可为任意 \(\lambda\in k\)。反过来，这些赋值满足唯一非零基本括号的关系，双线性保证在全部向量上成立，所以恰为全部一维表示。它们的核均包含 \(ke\ne0\)，因而没有一个忠实。
\end{solution}

\begin{optionalexercise}\label{b2:ex:20-heisenberg-adjoint-matrices}
对\cref{b2:exm:heisenberg-lie-algebra}中的 \(\mathfrak h\)，在有序基 \((x,d,z)\) 下求 \(\operatorname{ad}(x),\operatorname{ad}(d),\operatorname{ad}(z)\) 的矩阵及伴随表示的核。与该例中 \(\mathfrak h\subseteq M_3(k)^-\) 的自然表示比较。
\end{optionalexercise}
\begin{solution}
由 \([x,d]=-z\)、\([d,x]=z\) 及 \(z\) 居中，三矩阵依次为
\[
\begin{pmatrix}0&0&0\\0&0&0\\0&-1&0\end{pmatrix},\qquad
\begin{pmatrix}0&0&0\\0&0&0\\1&0&0\end{pmatrix},\qquad 0.
\]
前两个矩阵线性无关，所以核恰为 \(kz\)。这也由\cref{b2:prop:adjoint-lie-representation}等于中心：若 \(ax+bd+cz\) 居中，与 \(x\) 取括号得 \(bz=0\)，与 \(d\) 取括号得 \(-az=0\)，故 \(a=b=0\)。自然表示把 \(x,d,z\) 分别送到线性无关的 \(E_{23},E_{12},E_{13}\)，所以忠实。伴随表示不忠实，并不妨碍同一个 Lie 代数有别的忠实表示；这些计算在特征二时仍然成立。
\end{solution}

\begin{optionalexercise}\label{b2:ex:20-abelian-lie-jordan-modules}
令 \(\mathfrak a=ku\) 的括号为零，在 \(V=k^3\) 上规定 \(u\) 作用为
\[
J=\begin{pmatrix}0&1&0\\0&0&1\\0&0&0\end{pmatrix}.
\]
将这个表示对应的 \(U(\mathfrak a)\) 模写成一个循环模商，求 \(U(\mathfrak a)\) 对 \(V\) 的作用核，并求与 \(J\) 交换的全部线性算子。
\end{optionalexercise}
\begin{solution}
括号为零，且只有一个给定算子，所以确为表示。由\cref{b2:def:symmetric-algebra,b2:thm:symmetric-polynomial-algebra}及泛包络的关系，\(U(\mathfrak a)=k[t]\)，其中 \(t=j(u)\)。向量 \(v=e_3\) 满足 \(Jv=e_2,J^2v=e_1,J^3v=0\)，故映射 \(k[t]\to V\)、\(f\mapsto f(J)v\) 满射，其核为 \((t^3)\)：除以 \(t^3\) 所得次数小于三的余式，作用于 \(v\) 后的系数正是独立向量 \(e_3,e_2,e_1\) 的系数。因此 \(V\cong k[t]/(t^3)\)，作用核同样为 \((t^3)\)。

若 \(FJ=JF\)，则 \(F\) 由 \(F(v)\) 决定，因为 \(F(J^rv)=J^rF(v)\)。唯一写成 \(F(v)=(aI+bJ+cJ^2)v\) 后，这个等式在整组循环基上成立，故 \(F=aI+bJ+cJ^2\)。反过来，这些算子都与 \(J\) 交换。特别地，Lie 表示 \(ku\to\operatorname{End}_k(V)\) 本身单射，而其泛包络的作用不单射，二者的忠实性不可混同。
\end{solution}

\begin{optionalexercise}\label{b2:ex:20-enveloping-ideal-quotient}
设 \(I\) 是 Lie 代数 \(\mathfrak g\) 的理想，\(\bigl(j(I)\bigr)\) 表示 \(U(\mathfrak g)\) 中由 \(j(I)\) 生成的双边理想。不使用 PBW 定理，证明
\[
U(\mathfrak g/I)\cong U(\mathfrak g)/\bigl(j(I)\bigr).
\]
说明两边的模对应于哪些 \(\mathfrak g\) 表示。
\end{optionalexercise}
\begin{solution}
记 \(\pi:\mathfrak g\to\mathfrak g/I\) 为商 Lie 同态，\(j'\) 为商 Lie 代数到其泛包络的典范映射。\(j'\pi\) 由\cref{b2:thm:enveloping-universal}延拓为 \(U(\mathfrak g)\to U(\mathfrak g/I)\)，并零化 \(j(I)\)，故下降为同态 \(f:U(\mathfrak g)/\bigl(j(I)\bigr)\to U(\mathfrak g/I)\)。

反向令 \(q:U(\mathfrak g)\to U(\mathfrak g)/\bigl(j(I)\bigr)\) 为商同态。因 \(qj\) 零化 \(I\)，公式 \(u+I\mapsto qj(u)\) 不依赖代表元；\eqref{b2:eq:enveloping-canonical-lie-map}又保证它保持括号。再次应用泛性质，得到 \(g:U(\mathfrak g/I)\to U(\mathfrak g)/\bigl(j(I)\bigr)\)。两复合分别固定全部 \(qj(u)\) 与 \(j'(u+I)\)，这些元素连同单位元生成各自的代数，所以两复合为恒等，得到同构。

由\cref{b2:cor:lie-representations-enveloping-modules}，模掉 \(\bigl(j(I)\bigr)\) 正是要求 \(I\) 的每个元素在表示空间上作用为零。这也恰是 Lie 表示下降到 \(\mathfrak g/I\) 的条件。
\end{solution}

\begin{optionalexercise}\label{b2:ex:20-heisenberg-central-character}
对三维 Heisenberg Lie 代数 \(\mathfrak h\)，在特征零时构造一个三维忠实表示，使中心元素 \(z\) 作用非零；说明它为何不违背正文关于非零标量中心作用的有限维表示结论。特征 \(p>0\) 时，对任意 \(\lambda\in k^\times\) 构造一个 \(p\) 维简单表示，使 \(z\) 作用为 \(\lambda I\)，并证明这个 Lie 表示也忠实。
\end{optionalexercise}
\begin{solution}
在 \(k^3\) 上取 \(\rho(x)=E_{23},\rho(d)=E_{12},\rho(z)=E_{13}\)，则 \([E_{12},E_{23}]=E_{13}\)，且 \(E_{13}\) 与另两个矩阵交换。三个矩阵线性无关，故表示忠实；中心作用是非零平方零矩阵，并非非零标量矩阵。\cref{b2:prop:heisenberg-enveloping-weyl,b2:prop:weyl-finite-dimensional-representations}所排除的是特征零中中心作用为指定非零标量的情形，不能由此排除所有有限维 Heisenberg 表示。

特征 \(p\) 时，在 \(k[X]/(X^p)\) 上取 \(\rho(x)f=Xf\)、\(\rho(d)f=\lambda f'\)、\(\rho(z)f=\lambda f\)。关系成立，非零 \(\lambda\) 的缩放又不改变截断 Weyl 表示的子模，故它简单。若 \(a\rho(x)+b\rho(d)+c\rho(z)=0\)，作用于一得到 \(aX+c\lambda=0\)，从而 \(a=c=0\)；再作用于 \(X\) 得 \(b\lambda=0\)，故 \(b=0\)。因此三个作用算子独立，Lie 表示忠实。这里用到 \(p\ge2\)，所以一和 \(X\) 在商中确实独立。
\end{solution}

\begin{optionalexercise}\label{b2:ex:20-inner-sigma-derivation}
设 \(R\) 为含幺环，\(\sigma:R\to R\) 为保幺环自同态，固定 \(c\in R\)，令 \(\delta(a)=ca-\sigma(a)c\)。证明 \(\delta\) 是 \(\sigma\) 导子，并通过变量代换证明 \(R[t;\sigma,\delta]\cong R[u;\sigma]\)。
\end{optionalexercise}
\begin{solution}
直接展开得到
\[
\sigma(a)\delta(b)+\delta(a)b
=\sigma(a)cb-\sigma(a)\sigma(b)c+cab-\sigma(a)cb
=\delta(ab).
\]
在原扩张中置 \(u=t-c\)，则
\[
ua=\sigma(a)t+ca-\sigma(a)c-ca=\sigma(a)u.
\]
所以 Ore 泛性质给出从 \(R[u;\sigma]\) 到原扩张的同态。反向将 \(t\) 送到 \(u+c\)，同一个等式反向验证所需关系。两复合固定 \(R\) 和扩张变量，故互逆。
\end{solution}

\begin{optionalexercise}\label{b2:ex:20-ore-noninjective-sidedness}
令 \(R=k[x]\)、\(\sigma(f)=f(0)\)、\(S=R[t;\sigma]\)。对每个 \(m\ge1\)，分别求
\[
\{q\in S\mid t^mq=0\},\qquad \{q\in S\mid qt^m=0\}.
\]
证明第一个集合是双边理想，辨认相应商环，并说明这两个集合为何不能混同。
\end{optionalexercise}
\begin{solution}
把 \(q\) 唯一写成 \(\sum_i f_i(x)t^i\)。由于 \(\sigma^m(f)=f(0)\)，有 \(t^mq=\sum_i f_i(0)t^{m+i}\)，故它为零当且仅当每个 \(f_i\in xk[x]\)。另一方面，\(qt^m=\sum_i f_i(x)t^{m+i}\)，左正规式的独立性使其为零当且仅当 \(q=0\)。

第一个集合为 \(J=\bigoplus_{i\ge0}xk[x]t^i=xS\)。它显然对右乘封闭；左乘 \(k[x]\) 保持系数被 \(x\) 整除，左乘 \(t\) 则将全部元素杀掉，所以也对左乘封闭。将系数在零处求值、保留 \(t\)，得到满同态 \(S\to k[t]\)，核恰为 \(J\)，故 \(S/J\cong k[t]\)。非单射的 \(\sigma\) 能造成完全不同的左右零化行为；两侧集合的计算依靠左正规式，不可交换乘积次序。
\end{solution}

\begin{optionalexercise}\label{b2:ex:20-q-derivation}
固定 \(q\in k^\times\)，在 \(k[x]\) 上令 \(\sigma\bigl(f(x)\bigr)=f(qx)\)，并按线性定义
\[
\delta(x^n)=(1+q+\cdots+q^{n-1})x^{n-1}\quad(n\ge1),
\qquad\delta(1)=0.
\]
证明这是左 \(\sigma\) 导子，并写出相应 Ore 扩张的 \(tx\)、\(tx^2\) 及正规单项式基。
\end{optionalexercise}
\begin{solution}
记 \([n]_q=1+q+\cdots+q^{n-1}\)，则 \([m+n]_q=[m]_q+q^m[n]_q\)。因此在两个单项式上，\(\delta(x^mx^n)=\sigma(x^m)\delta(x^n)+\delta(x^m)x^n\)。两边对多项式分别线性，故等式对任意乘积成立。由构造定理，
\[
tx=qxt+1,\qquad tx^2=q^2x^2t+(1+q)x,
\]
正规单项式基为 \(x^it^j\)。当 \(q=1\) 时，\([n]_q=n\)，恢复通常求导及 Weyl 关系；所有定义都没有使用除以 \(q-1\)。
\end{solution}

\begin{optionalexercise}\label{b2:ex:20-differential-line}
在微分域 \((K,\delta)\) 上取 \(L=\partial-a\)，写出 \(M_L\) 上的作用。若将基向量 \(e\) 改为 \(f=be\)、\(b\ne0\)，新作用系数是什么？最后在 \(K=k(x)\)、\(\delta=\mathrm{d}/\mathrm{d}x\)、\(a=1/x\) 时，用常数域 \(C=\ker\delta\) 表示全部 \(K\) 内的解。
\end{optionalexercise}
\begin{solution}
模以 \(e=1+DL\) 为 \(K\) 基，满足 \(\partial e=ae\)。因此 \(\partial(ue)=\bigl(\delta(u)+au\bigr)e\)。换基后
\[
\partial f=\bigl(\delta(b)+ab\bigr)e
=\bigl(a+b^{-1}\cdot\delta(b)\bigr)f.
\]
对于最后的方程，\(\delta(y)=y/x\) 等价于 \(\delta(y/x)=0\)，所以所有解为 \(y=cx\)、\(c\in C\)。这同时表明解空间自然只需对常数域封闭；在正特征中，\(C\) 可以大于给定的 \(k\)。
\end{solution}

\begin{optionalexercise}\label{b2:ex:20-differential-companion}
设 \((K,\delta)\) 为微分域，左微分模在一组基 \(e=(e_0,\ldots,e_{m-1})\) 中的坐标作用为 \(\nabla\mathbf v=\delta\mathbf v+A\mathbf v\)。令 \(e'=eP\)、\(P\in\operatorname{GL}_m(K)\)，求新作用矩阵和模同态取值列 \(Y'\) 的变换式。在 \(K=k(x)\)、\(L=\partial^2-x\) 时，取
\[
P=\begin{pmatrix}1&x\\0&1\end{pmatrix},
\]
实际计算这两个变换。
\end{optionalexercise}
\begin{solution}
同一元素满足 \(e\mathbf v=e'\mathbf v'\)，即 \(\mathbf v=P\mathbf v'\)。由 Leibniz 公式，
\[
\delta(P\mathbf v')+AP\mathbf v'
=P\delta\mathbf v'+(\delta P+AP)\mathbf v'.
\]
改为新基坐标，得 \(B=P^{-1}AP+P^{-1}\delta P\)。系数求导项不可丢弃。若 \(f\) 是到函数模的同态，\(Y_i=f(e_i)\)，则 \(Y'=P^{\mathsf T}Y\)。正文的系统 \(\delta Y=A^{\mathsf T}Y\) 因而变为 \(\delta Y'=B^{\mathsf T}Y'\)，这也可直接对 \(P^{\mathsf T}Y\) 求导验证。

对于 \(L=\partial^2-x\)，\cref{b2:prop:differential-companion-coordinates}给出 \(A=\begin{psmallmatrix}0&x\\1&0\end{psmallmatrix}\)。代入可得
\[
B=\begin{pmatrix}-x&x+1-x^2\\1&x\end{pmatrix},
\qquad Y'=\begin{pmatrix}y\\xy+\delta y\end{pmatrix}.
\]
对第二个分量求导，用 \(\delta^2y=xy\)，即得 \(\delta Y'=B^{\mathsf T}Y'\)。这里两个基之间的 \(K\) 线性换坐标并不将 \(\nabla\) 变成 \(K\) 线性算子。
\end{solution}

\begin{optionalexercise}\label{b2:ex:20-diamond-idempotent-rules}
用规则 \(x^2\to x\)、\(yx\to y\) 检查全部非平凡歧义，证明代数 \(k\langle x,y\rangle/(x^2-x,yx-y)\) 有基 \(y^n,xy^n\)、\(n\ge0\)。给出其中两个非零元素的乘积为零的例子。
\end{optionalexercise}
\begin{solution}
两条规则都降低字长。只有 \(xxx\) 与 \(yxx\) 存在重叠：\(xxx\) 无论先改哪一对都化为 \(x\)，\(yxx\) 两种改法都化为 \(y\)。没有非平凡包含歧义，所以\cref{b2:thm:diamond-lemma}适用。不可约字中不能有连续的两个 \(x\)，也不能有 \(y\) 在 \(x\) 左边，故恰为所列形式。正规字基保证 \(y\ne0\)、\(1-x\ne0\)，但关系给出 \(y(1-x)=0\)。
\end{solution}

\begin{optionalexercise}\label{b2:ex:20-diamond-two-failures}
固定 \(a,b,c\in k\)，在 \(k\langle x,y\rangle\) 中取规则 \(x^2\to ax+b\)、\(yx\to cy\)。证明这组规则的全部歧义可解当且仅当 \(c^2-ac-b=0\)。在该条件成立时给出商代数的一组基；条件不成立时，辨认隐藏关系及实际商代数。

再对 \(\alpha,\beta\in k\) 考察规则 \(xyx\to\alpha x\)、\(yx\to\beta\) 的包含歧义，辨认相应商代数，并在非零商的情形给出一组基。
\end{optionalexercise}
\begin{solution}
两条规则降低字长，非平凡重叠只有 \(xxx\) 与 \(yxx\)，没有非平凡包含。前者两支都成为 \((a^2+b)x+ab\)，后者两支分别为 \(c^2y\) 与 \((ac+b)y\)。它们已是正规式，所以歧义可解恰好要求 \(c^2-ac-b=0\)。在该条件下，Diamond 引理给出不可约字基 \(y^n,xy^n\)、\(n\ge0\)：每个 \(x\) 只能位于字首，且至多有一个。

若 \(c^2-ac-b\ne0\)，重叠差强迫 \(y=0\)。将 \(y\) 送到零、\(x\) 送到自身，得到商到 \(k[x]/(x^2-ax-b)\) 的同态；反向保留 \(x\) 的类，利用 \(y=0\)，两复合为恒等。因此实际商正是这个二次多项式商，基为 \(1,x\)。这既计算出额外关系，也表明规则不相容时不能沿用原先预期的无限正规字基。

第二组规则在 \(xyx\) 处，整体改写得到 \(\alpha x\)，改写其中的 \(yx\) 得到 \(\beta x\)。若 \(\alpha=\beta\)，第一条关系由第二条左乘 \(x\) 得到，故商只需关系 \(yx=\beta\)。规则 \(yx\to\beta\) 没有自身包含或重叠，Diamond 引理给出基 \(x^iy^j\)、\(i,j\ge0\)。若 \(\alpha\ne\beta\)，包含差强迫 \(x=0\)，第二条关系随即要求 \(\beta1=0\)。当 \(\beta=0\) 时，正反赋值 \(x\mapsto0\)、保留 \(y\) 给出商 \(k[y]\)，其基为 \(y^j\)；当 \(\beta\ne0\) 时，\(1=0\)，商为零代数。这一包含歧义可以删掉一个生成元，也可以使整个商为零。
\end{solution}

\begin{optionalexercise}\label{b2:ex:20-infinite-groebner-growth}
对 \(B=k\langle x,y\rangle/(x^2-xy)\) 取字长分次。利用\cref{b2:exm:infinite-noncommutative-groebner}求 \(\dim_kB_n\) 及形式 Hilbert 级数 \(\sum_{n\ge0}\dim_k(B_n)t^n\)，并证明 \(B\) 虽与二元多项式环具有相同的分次维数，却有零因子。
\end{optionalexercise}
\begin{solution}
次数 \(n\ge1\) 的正规字中，有一个 \(y^n\)，以及满足 \(a+b=n-1\) 的 \(n\) 个 \(y^axy^b\)，所以维数为 \(n+1\)；零次维数为一。因此级数为
\[
\sum_{n\ge0}(n+1)t^n=\frac1{(1-t)^2},
\]
这里是形式幂级数恒等式。正规字基使 \(x\) 与 \(x-y\) 均非零，但 \(x(x-y)=0\)。所以分次维数不能决定乘法结构，也不能单凭这一数列判断无零因子性。
\end{solution}

\begin{optionalexercise}\label{b2:ex:20-groebner-order-dependence}
在 \(B=k\langle x,y\rangle/(x^2-xy)\) 中，使用\cref{b2:exm:infinite-noncommutative-groebner}关于 \(y>x\) 的正规基，计算 \((y^ax^b)(y^cx^d)\)，其中 \(a,b,c,d\ge0\)。据此求中心，并证明中心为 \(k\) 并不使 \(B\) 成为单环。
\end{optionalexercise}
\begin{solution}
由 \(xy=x^2\)，归纳得 \(x^by^c=x^{b+c}\) 对 \(b\ge1\) 成立。因此
\[
(y^ax^b)(y^cx^d)=
\begin{cases}
y^{a+c}x^d,&b=0,\\
y^ax^{b+c+d},&b\ge1.
\end{cases}
\]
若 \(z=\sum_{b=0}^Nf_b(y)x^b\) 与 \(y\) 交换，且 \(N\ge1,f_N\ne0\)，则 \(zy\) 的 \(x^{N+1}\) 系数为 \(f_N(y)\)，\(yz\) 却没有这一项，违反正规基的独立性。因此 \(z=f_0(y)\)。再要求它与 \(x\) 交换：若 \(f_0\) 的次数 \(n\ge1\)，\(xf_0(y)\) 含非零 \(x^{n+1}\) 项，而 \(f_0(y)x\) 只有 \(x\) 次数一，仍矛盾。故 \(Z(B)=k\)。

双边理想 \((x)\) 非零，因为 \(x\) 是正规基元素；它又是真理想，因为模掉 \(x\) 得到非零商 \(k[y]\)。所以 \(B\) 不是单环。不同字序下的正规式允许不同计算，但所得中心和理想仍是同一个代数的内在对象。
\end{solution}

\begin{optionalexercise}\label{b2:ex:20-differential-associated-module}
在 \(D=K[\partial;\delta]\) 中按 \(\partial\) 次数滤过，\(K\) 在零次。若 \(L\) 为 \(m\) 阶首一算子，证明 \(\operatorname{gr}D\cong K[T]\)，并对商的诱导滤过证明
\[
\operatorname{gr}(D/DL)\cong K[T]/(T^m)
\]
为分次左模。
\end{optionalexercise}
\begin{solution}
关系 \(\partial a-a\partial=\delta(a)\) 降低一次 \(\partial\) 次数，故最高次符号与 \(K\) 交换；Ore 正规式使每个次数商片一维，得到 \(\operatorname{gr}D\cong K[T]\)。对于左理想 \(DL\)，非零 \(QL\) 的次数为 \(\deg Q+m\)，其主符号为 \(\operatorname{symb}(Q)T^m\)，所以该理想的伴随分次恰为 \((T^m)\)。也可直接使用唯一左余式：商模在次数小于 \(m\) 时各新增一个基向量，随后不再新增，且乘 \(T\) 逐次移位，最后为零。于是得到所写的分次模同构。
\end{solution}

\begin{optionalexercise}\label{b2:ex:20-quantum-matrix-fibre}
设 \(\operatorname{char}k\ne2\)，\(a=s^2\ne0\)、\(b\ne0\) 为 \(k\) 中元素。证明
\[
k_{-1}[x,y]/(x^2-a,y^2-b)\cong M_2(k).
\]
\end{optionalexercise}
\begin{solution}
由量子平面的正规单项式基及平方关系，商由 \(1,x,y,xy\) 线性生成。取
\[
X=\begin{pmatrix}s&0\\0&-s\end{pmatrix},\qquad
Y=\begin{pmatrix}0&b\\1&0\end{pmatrix}.
\]
它们满足 \(YX=-XY\)、\(X^2=aI\)、\(Y^2=bI\)，所以给出到矩阵代数的同态。\(I,X\) 线性独立并生成对角矩阵；\(Y,XY\) 的两个非对角位置组成的系数矩阵行列式为 \(-2sb\ne0\)，所以生成全部非对角矩阵。因此像为四维矩阵代数，来源又至多四维，所得满同态必为同构。
\end{solution}

\begin{optionalexercise}\label{b2:ex:20-ore-right-normal-form}
设 \(R\) 为含幺环，\(\sigma\) 为环自同构，\(\delta\) 为左 \(\sigma\) 导子。置 \(\tau=\sigma^{-1}\)、\(\varepsilon=-\delta\sigma^{-1}\)。把 \(\tau,\varepsilon\) 视为 \(R^{\mathrm{op}}\) 上的映射，证明 \(\varepsilon\) 是左 \(\tau\) 导子，并证明
\[
R[t;\sigma,\delta]^{\mathrm{op}}\cong R^{\mathrm{op}}[u;\tau,\varepsilon].
\]
同构在底环上为恒等，将 \(u\) 送到反环中的 \(t\)。
\end{optionalexercise}
\begin{solution}
在原环记乘积，扭曲 Leibniz 公式给出
\[
\varepsilon(ba)=b\varepsilon(a)+\varepsilon(b)\tau(a).
\]
在反环中它恰好是 \(\varepsilon(a\cdot_{\mathrm{op}}b)=\tau(a)\cdot_{\mathrm{op}}\varepsilon(b)+\varepsilon(a)\cdot_{\mathrm{op}}b\)，故是所需导子。\(\tau\) 也是反环自同构。

在原 Ore 扩张中，\(at=t\tau(a)+\varepsilon(a)\)。反过来读乘积即为 \(t\cdot_{\mathrm{op}}a=\tau(a)\cdot_{\mathrm{op}}t+\varepsilon(a)\)。Ore 泛性质给出题设同态；它的像包含底环和 \(t\)，故满射。由\cref{b2:prop:ore-right-normal-form}，原扩张以 \(t^i\) 为右 \(R\) 模基，因而其反环以同一组元素为左 \(R^{\mathrm{op}}\) 模基。这个同态将 Ore 左正规基逐项送到该基，所以单射。扭曲导子中的负号和逆自同构由反转乘法次序产生。
\end{solution}
